\documentclass[11pt]{amsart}
\usepackage[margin=1.1in]{geometry}
\usepackage{amssymb,mathtools,booktabs,multicol}
\usepackage[hidelinks]{hyperref}
\newcommand{\bclosed}{243}
\newcommand{\bopen}{293}
\newcommand{\btotal}{536}
\newcommand{\bgrouptable}{$\O_4(\F_{5})$ & reflections & 15 \\
$\Sp_4(\F_{3})$ & transvections & 67 \\
$\Sp_4(\F_{5})$ & transvections & 161 \\}
\newcommand{\bclosedlist}{$13\mathrm{a}_{69}$, $13\mathrm{a}_{70}$, $13\mathrm{a}_{72}$, $13\mathrm{a}_{73}$, $13\mathrm{a}_{74}$, $13\mathrm{a}_{76}$, $13\mathrm{a}_{79}$, $13\mathrm{a}_{96}$, $13\mathrm{a}_{104}$, $13\mathrm{a}_{106}$, $13\mathrm{a}_{107}$, $13\mathrm{a}_{108}$, $13\mathrm{a}_{156}$, $13\mathrm{a}_{157}$, $13\mathrm{a}_{160}$, $13\mathrm{a}_{171}$, $13\mathrm{a}_{235}$, $13\mathrm{a}_{343}$, $13\mathrm{a}_{353}$, $13\mathrm{a}_{354}$, $13\mathrm{a}_{355}$, $13\mathrm{a}_{358}$, $13\mathrm{a}_{359}$, $13\mathrm{a}_{360}$, $13\mathrm{a}_{388}$, $13\mathrm{a}_{390}$, $13\mathrm{a}_{396}$, $13\mathrm{a}_{399}$, $13\mathrm{a}_{400}$, $13\mathrm{a}_{402}$, $13\mathrm{a}_{403}$, $13\mathrm{a}_{408}$, $13\mathrm{a}_{413}$, $13\mathrm{a}_{414}$, $13\mathrm{a}_{505}$, $13\mathrm{a}_{524}$, $13\mathrm{a}_{525}$, $13\mathrm{a}_{526}$, $13\mathrm{a}_{527}$, $13\mathrm{a}_{528}$, $13\mathrm{a}_{529}$, $13\mathrm{a}_{535}$, $13\mathrm{a}_{536}$, $13\mathrm{a}_{541}$, $13\mathrm{a}_{542}$, $13\mathrm{a}_{543}$, $13\mathrm{a}_{544}$, $13\mathrm{a}_{545}$, $13\mathrm{a}_{554}$, $13\mathrm{a}_{555}$, $13\mathrm{a}_{557}$, $13\mathrm{a}_{558}$, $13\mathrm{a}_{596}$, $13\mathrm{a}_{597}$, $13\mathrm{a}_{601}$, $13\mathrm{a}_{621}$, $13\mathrm{a}_{630}$, $13\mathrm{a}_{631}$, $13\mathrm{a}_{632}$, $13\mathrm{a}_{636}$, $13\mathrm{a}_{637}$, $13\mathrm{a}_{657}$, $13\mathrm{a}_{664}$, $13\mathrm{a}_{665}$, $13\mathrm{a}_{815}$, $13\mathrm{a}_{867}$, $13\mathrm{a}_{1392}$, $13\mathrm{a}_{1538}$, $13\mathrm{a}_{1995}$, $13\mathrm{a}_{2006}$, $13\mathrm{a}_{2008}$, $13\mathrm{a}_{2613}$, $13\mathrm{a}_{2614}$, $13\mathrm{a}_{2615}$, $13\mathrm{a}_{2616}$, $13\mathrm{a}_{2617}$, $13\mathrm{a}_{2619}$, $13\mathrm{a}_{2620}$, $13\mathrm{a}_{2800}$, $13\mathrm{a}_{2801}$, $13\mathrm{a}_{3038}$, $13\mathrm{a}_{3040}$, $13\mathrm{a}_{3107}$, $13\mathrm{a}_{3144}$, $13\mathrm{a}_{3146}$, $13\mathrm{a}_{3179}$, $13\mathrm{a}_{3257}$, $13\mathrm{a}_{3333}$, $13\mathrm{a}_{3334}$, $13\mathrm{a}_{3343}$, $13\mathrm{a}_{3344}$, $13\mathrm{a}_{3376}$, $13\mathrm{a}_{3377}$, $13\mathrm{a}_{3378}$, $13\mathrm{a}_{3380}$, $13\mathrm{a}_{3381}$, $13\mathrm{n}_{160}$, $13\mathrm{n}_{161}$, $13\mathrm{n}_{162}$, $13\mathrm{n}_{163}$, $13\mathrm{n}_{164}$, $13\mathrm{n}_{165}$, $13\mathrm{n}_{166}$, $13\mathrm{n}_{167}$, $13\mathrm{n}_{168}$, $13\mathrm{n}_{169}$, $13\mathrm{n}_{173}$, $13\mathrm{n}_{174}$, $13\mathrm{n}_{175}$, $13\mathrm{n}_{176}$, $13\mathrm{n}_{177}$, $13\mathrm{n}_{178}$, $13\mathrm{n}_{179}$, $13\mathrm{n}_{180}$, $13\mathrm{n}_{181}$, $13\mathrm{n}_{182}$, $13\mathrm{n}_{211}$, $13\mathrm{n}_{429}$, $13\mathrm{n}_{430}$, $13\mathrm{n}_{431}$, $13\mathrm{n}_{432}$, $13\mathrm{n}_{433}$, $13\mathrm{n}_{434}$, $13\mathrm{n}_{435}$, $13\mathrm{n}_{440}$, $13\mathrm{n}_{441}$, $13\mathrm{n}_{442}$, $13\mathrm{n}_{449}$, $13\mathrm{n}_{467}$, $13\mathrm{n}_{601}$, $13\mathrm{n}_{609}$, $13\mathrm{n}_{610}$, $13\mathrm{n}_{929}$, $13\mathrm{n}_{930}$, $13\mathrm{n}_{931}$, $13\mathrm{n}_{932}$, $13\mathrm{n}_{933}$, $13\mathrm{n}_{957}$, $13\mathrm{n}_{958}$, $13\mathrm{n}_{959}$, $13\mathrm{n}_{960}$, $13\mathrm{n}_{961}$, $13\mathrm{n}_{962}$, $13\mathrm{n}_{963}$, $13\mathrm{n}_{964}$, $13\mathrm{n}_{965}$, $13\mathrm{n}_{966}$, $13\mathrm{n}_{967}$, $13\mathrm{n}_{968}$, $13\mathrm{n}_{977}$, $13\mathrm{n}_{978}$, $13\mathrm{n}_{979}$, $13\mathrm{n}_{980}$, $13\mathrm{n}_{981}$, $13\mathrm{n}_{982}$, $13\mathrm{n}_{983}$, $13\mathrm{n}_{1042}$, $13\mathrm{n}_{1058}$, $13\mathrm{n}_{1059}$, $13\mathrm{n}_{1060}$, $13\mathrm{n}_{1061}$, $13\mathrm{n}_{1062}$, $13\mathrm{n}_{1063}$, $13\mathrm{n}_{1064}$, $13\mathrm{n}_{1072}$, $13\mathrm{n}_{1073}$, $13\mathrm{n}_{1074}$, $13\mathrm{n}_{1075}$, $13\mathrm{n}_{1076}$, $13\mathrm{n}_{1077}$, $13\mathrm{n}_{1078}$, $13\mathrm{n}_{1082}$, $13\mathrm{n}_{1083}$, $13\mathrm{n}_{1084}$, $13\mathrm{n}_{1103}$, $13\mathrm{n}_{1104}$, $13\mathrm{n}_{1105}$, $13\mathrm{n}_{1106}$, $13\mathrm{n}_{1239}$, $13\mathrm{n}_{1280}$, $13\mathrm{n}_{1281}$, $13\mathrm{n}_{1282}$, $13\mathrm{n}_{1283}$, $13\mathrm{n}_{1284}$, $13\mathrm{n}_{1285}$, $13\mathrm{n}_{1286}$, $13\mathrm{n}_{1287}$, $13\mathrm{n}_{1288}$, $13\mathrm{n}_{1289}$, $13\mathrm{n}_{1290}$, $13\mathrm{n}_{1291}$, $13\mathrm{n}_{1292}$, $13\mathrm{n}_{1293}$, $13\mathrm{n}_{1294}$, $13\mathrm{n}_{1295}$, $13\mathrm{n}_{1305}$, $13\mathrm{n}_{1306}$, $13\mathrm{n}_{1307}$, $13\mathrm{n}_{1308}$, $13\mathrm{n}_{1309}$, $13\mathrm{n}_{1310}$, $13\mathrm{n}_{1311}$, $13\mathrm{n}_{1312}$, $13\mathrm{n}_{1313}$, $13\mathrm{n}_{1314}$, $13\mathrm{n}_{1315}$, $13\mathrm{n}_{1316}$, $13\mathrm{n}_{1317}$, $13\mathrm{n}_{1318}$, $13\mathrm{n}_{1319}$, $13\mathrm{n}_{1320}$, $13\mathrm{n}_{1321}$, $13\mathrm{n}_{1322}$, $13\mathrm{n}_{1323}$, $13\mathrm{n}_{1324}$, $13\mathrm{n}_{1332}$, $13\mathrm{n}_{1333}$, $13\mathrm{n}_{1334}$, $13\mathrm{n}_{1363}$, $13\mathrm{n}_{1377}$, $13\mathrm{n}_{1378}$, $13\mathrm{n}_{1379}$, $13\mathrm{n}_{1380}$, $13\mathrm{n}_{1381}$, $13\mathrm{n}_{1385}$, $13\mathrm{n}_{1386}$, $13\mathrm{n}_{2933}$, $13\mathrm{n}_{2934}$, $13\mathrm{n}_{2935}$, $13\mathrm{n}_{2936}$, $13\mathrm{n}_{2937}$, $13\mathrm{n}_{2954}$, $13\mathrm{n}_{2955}$, $13\mathrm{n}_{2956}$, $13\mathrm{n}_{3390}$, $13\mathrm{n}_{3505}$, $13\mathrm{n}_{3723}$, $13\mathrm{n}_{3725}$, $13\mathrm{n}_{3726}$, $13\mathrm{n}_{3727}$, $13\mathrm{n}_{3763}$, $13\mathrm{n}_{3838}$, $13\mathrm{n}_{3915}$}
\newcommand{\bopenlist}{$13\mathrm{a}_{129}$, $13\mathrm{a}_{135}$, $13\mathrm{a}_{136}$, $13\mathrm{a}_{137}$, $13\mathrm{a}_{194}$, $13\mathrm{a}_{196}$, $13\mathrm{a}_{197}$, $13\mathrm{a}_{198}$, $13\mathrm{a}_{233}$, $13\mathrm{a}_{234}$, $13\mathrm{a}_{239}$, $13\mathrm{a}_{242}$, $13\mathrm{a}_{243}$, $13\mathrm{a}_{244}$, $13\mathrm{a}_{253}$, $13\mathrm{a}_{265}$, $13\mathrm{a}_{268}$, $13\mathrm{a}_{269}$, $13\mathrm{a}_{270}$, $13\mathrm{a}_{271}$, $13\mathrm{a}_{272}$, $13\mathrm{a}_{304}$, $13\mathrm{a}_{309}$, $13\mathrm{a}_{338}$, $13\mathrm{a}_{362}$, $13\mathrm{a}_{365}$, $13\mathrm{a}_{366}$, $13\mathrm{a}_{369}$, $13\mathrm{a}_{373}$, $13\mathrm{a}_{375}$, $13\mathrm{a}_{387}$, $13\mathrm{a}_{391}$, $13\mathrm{a}_{392}$, $13\mathrm{a}_{471}$, $13\mathrm{a}_{483}$, $13\mathrm{a}_{487}$, $13\mathrm{a}_{516}$, $13\mathrm{a}_{518}$, $13\mathrm{a}_{530}$, $13\mathrm{a}_{623}$, $13\mathrm{a}_{629}$, $13\mathrm{a}_{3145}$, $13\mathrm{a}_{3147}$, $13\mathrm{a}_{3148}$, $13\mathrm{a}_{3149}$, $13\mathrm{a}_{3167}$, $13\mathrm{a}_{3168}$, $13\mathrm{a}_{3169}$, $13\mathrm{a}_{3170}$, $13\mathrm{a}_{3188}$, $13\mathrm{a}_{3189}$, $13\mathrm{a}_{3190}$, $13\mathrm{a}_{3191}$, $13\mathrm{a}_{3204}$, $13\mathrm{a}_{3205}$, $13\mathrm{a}_{3207}$, $13\mathrm{a}_{3208}$, $13\mathrm{a}_{3209}$, $13\mathrm{a}_{3210}$, $13\mathrm{a}_{3213}$, $13\mathrm{a}_{3216}$, $13\mathrm{a}_{3219}$, $13\mathrm{a}_{3220}$, $13\mathrm{a}_{3221}$, $13\mathrm{a}_{3222}$, $13\mathrm{a}_{3223}$, $13\mathrm{a}_{3232}$, $13\mathrm{a}_{3234}$, $13\mathrm{a}_{3248}$, $13\mathrm{a}_{3259}$, $13\mathrm{a}_{3261}$, $13\mathrm{a}_{3262}$, $13\mathrm{a}_{3263}$, $13\mathrm{a}_{3264}$, $13\mathrm{a}_{3266}$, $13\mathrm{a}_{3274}$, $13\mathrm{a}_{3275}$, $13\mathrm{a}_{3276}$, $13\mathrm{a}_{3291}$, $13\mathrm{a}_{3292}$, $13\mathrm{a}_{3293}$, $13\mathrm{a}_{3294}$, $13\mathrm{a}_{3295}$, $13\mathrm{a}_{3298}$, $13\mathrm{a}_{3301}$, $13\mathrm{a}_{3302}$, $13\mathrm{a}_{3303}$, $13\mathrm{a}_{3317}$, $13\mathrm{a}_{3320}$, $13\mathrm{a}_{3323}$, $13\mathrm{a}_{3330}$, $13\mathrm{a}_{3331}$, $13\mathrm{a}_{3335}$, $13\mathrm{a}_{3345}$, $13\mathrm{a}_{3346}$, $13\mathrm{a}_{3347}$, $13\mathrm{a}_{3348}$, $13\mathrm{a}_{3351}$, $13\mathrm{a}_{3354}$, $13\mathrm{a}_{3355}$, $13\mathrm{a}_{3356}$, $13\mathrm{a}_{3363}$, $13\mathrm{a}_{3365}$, $13\mathrm{a}_{3368}$, $13\mathrm{a}_{3374}$, $13\mathrm{a}_{3375}$, $13\mathrm{a}_{3388}$, $13\mathrm{a}_{3390}$, $13\mathrm{a}_{3393}$, $13\mathrm{a}_{3394}$, $13\mathrm{a}_{3395}$, $13\mathrm{a}_{3396}$, $13\mathrm{a}_{3399}$, $13\mathrm{a}_{3400}$, $13\mathrm{a}_{3401}$, $13\mathrm{a}_{3402}$, $13\mathrm{a}_{3403}$, $13\mathrm{a}_{3404}$, $13\mathrm{a}_{3405}$, $13\mathrm{a}_{3407}$, $13\mathrm{a}_{3409}$, $13\mathrm{a}_{3410}$, $13\mathrm{a}_{3411}$, $13\mathrm{a}_{3412}$, $13\mathrm{a}_{3413}$, $13\mathrm{a}_{3415}$, $13\mathrm{a}_{3417}$, $13\mathrm{a}_{3418}$, $13\mathrm{a}_{3811}$, $13\mathrm{a}_{3812}$, $13\mathrm{n}_{326}$, $13\mathrm{n}_{329}$, $13\mathrm{n}_{330}$, $13\mathrm{n}_{341}$, $13\mathrm{n}_{344}$, $13\mathrm{n}_{345}$, $13\mathrm{n}_{348}$, $13\mathrm{n}_{351}$, $13\mathrm{n}_{352}$, $13\mathrm{n}_{355}$, $13\mathrm{n}_{357}$, $13\mathrm{n}_{358}$, $13\mathrm{n}_{510}$, $13\mathrm{n}_{513}$, $13\mathrm{n}_{514}$, $13\mathrm{n}_{518}$, $13\mathrm{n}_{520}$, $13\mathrm{n}_{521}$, $13\mathrm{n}_{524}$, $13\mathrm{n}_{525}$, $13\mathrm{n}_{526}$, $13\mathrm{n}_{529}$, $13\mathrm{n}_{530}$, $13\mathrm{n}_{531}$, $13\mathrm{n}_{582}$, $13\mathrm{n}_{590}$, $13\mathrm{n}_{591}$, $13\mathrm{n}_{598}$, $13\mathrm{n}_{599}$, $13\mathrm{n}_{600}$, $13\mathrm{n}_{627}$, $13\mathrm{n}_{632}$, $13\mathrm{n}_{633}$, $13\mathrm{n}_{644}$, $13\mathrm{n}_{647}$, $13\mathrm{n}_{648}$, $13\mathrm{n}_{651}$, $13\mathrm{n}_{652}$, $13\mathrm{n}_{661}$, $13\mathrm{n}_{662}$, $13\mathrm{n}_{663}$, $13\mathrm{n}_{694}$, $13\mathrm{n}_{695}$, $13\mathrm{n}_{696}$, $13\mathrm{n}_{704}$, $13\mathrm{n}_{705}$, $13\mathrm{n}_{706}$, $13\mathrm{n}_{707}$, $13\mathrm{n}_{708}$, $13\mathrm{n}_{709}$, $13\mathrm{n}_{712}$, $13\mathrm{n}_{713}$, $13\mathrm{n}_{716}$, $13\mathrm{n}_{717}$, $13\mathrm{n}_{718}$, $13\mathrm{n}_{834}$, $13\mathrm{n}_{838}$, $13\mathrm{n}_{839}$, $13\mathrm{n}_{857}$, $13\mathrm{n}_{926}$, $13\mathrm{n}_{991}$, $13\mathrm{n}_{992}$, $13\mathrm{n}_{993}$, $13\mathrm{n}_{1000}$, $13\mathrm{n}_{1001}$, $13\mathrm{n}_{1004}$, $13\mathrm{n}_{1005}$, $13\mathrm{n}_{1006}$, $13\mathrm{n}_{1007}$, $13\mathrm{n}_{1011}$, $13\mathrm{n}_{1012}$, $13\mathrm{n}_{1013}$, $13\mathrm{n}_{1039}$, $13\mathrm{n}_{1040}$, $13\mathrm{n}_{1041}$, $13\mathrm{n}_{1043}$, $13\mathrm{n}_{1044}$, $13\mathrm{n}_{1045}$, $13\mathrm{n}_{1046}$, $13\mathrm{n}_{1047}$, $13\mathrm{n}_{1048}$, $13\mathrm{n}_{1175}$, $13\mathrm{n}_{1178}$, $13\mathrm{n}_{1179}$, $13\mathrm{n}_{1200}$, $13\mathrm{n}_{1201}$, $13\mathrm{n}_{1202}$, $13\mathrm{n}_{1210}$, $13\mathrm{n}_{1211}$, $13\mathrm{n}_{1212}$, $13\mathrm{n}_{1270}$, $13\mathrm{n}_{1276}$, $13\mathrm{n}_{1296}$, $13\mathrm{n}_{1367}$, $13\mathrm{n}_{1370}$, $13\mathrm{n}_{1371}$, $13\mathrm{n}_{1374}$, $13\mathrm{n}_{1375}$, $13\mathrm{n}_{1376}$, $13\mathrm{n}_{3739}$, $13\mathrm{n}_{3742}$, $13\mathrm{n}_{3745}$, $13\mathrm{n}_{3748}$, $13\mathrm{n}_{3770}$, $13\mathrm{n}_{3773}$, $13\mathrm{n}_{3776}$, $13\mathrm{n}_{3779}$, $13\mathrm{n}_{3814}$, $13\mathrm{n}_{3840}$, $13\mathrm{n}_{3842}$, $13\mathrm{n}_{3844}$, $13\mathrm{n}_{3854}$, $13\mathrm{n}_{3855}$, $13\mathrm{n}_{3856}$, $13\mathrm{n}_{3870}$, $13\mathrm{n}_{3871}$, $13\mathrm{n}_{3872}$, $13\mathrm{n}_{3873}$, $13\mathrm{n}_{3874}$, $13\mathrm{n}_{3897}$, $13\mathrm{n}_{3898}$, $13\mathrm{n}_{3899}$, $13\mathrm{n}_{3900}$, $13\mathrm{n}_{3901}$, $13\mathrm{n}_{3903}$, $13\mathrm{n}_{3904}$, $13\mathrm{n}_{3905}$, $13\mathrm{n}_{3911}$, $13\mathrm{n}_{3914}$, $13\mathrm{n}_{3916}$, $13\mathrm{n}_{3918}$, $13\mathrm{n}_{3919}$, $13\mathrm{n}_{3920}$, $13\mathrm{n}_{3921}$, $13\mathrm{n}_{3922}$, $13\mathrm{n}_{3924}$, $13\mathrm{n}_{3925}$, $13\mathrm{n}_{3928}$, $13\mathrm{n}_{3929}$, $13\mathrm{n}_{3931}$, $13\mathrm{n}_{3932}$, $13\mathrm{n}_{3933}$, $13\mathrm{n}_{3934}$, $13\mathrm{n}_{3935}$, $13\mathrm{n}_{3936}$, $13\mathrm{n}_{3937}$, $13\mathrm{n}_{3938}$, $13\mathrm{n}_{3939}$, $13\mathrm{n}_{3940}$, $13\mathrm{n}_{3941}$, $13\mathrm{n}_{4587}$, $13\mathrm{n}_{4589}$, $13\mathrm{n}_{4633}$, $13\mathrm{n}_{4634}$, $13\mathrm{n}_{4637}$, $13\mathrm{n}_{4638}$, $13\mathrm{n}_{4639}$, $13\mathrm{n}_{4640}$, $13\mathrm{n}_{4641}$, $13\mathrm{n}_{4642}$, $13\mathrm{n}_{4643}$, $13\mathrm{n}_{4644}$, $13\mathrm{n}_{4645}$}

\newcommand{\ntotal}{3525}
\newcommand{\ncerts}{2728}
\newcommand{\nclosed}{2177}
\newcommand{\nclosedone}{520}
\newcommand{\nclosedtwo}{1657}
\newcommand{\nimproved}{551}
\newcommand{\nchained}{2}
\newcommand{\nlower}{382}
\newcommand{\nlowertwo}{360}
\newcommand{\nlowerthree}{22}
\newcommand{\nlowerten}{314}
\newcommand{\nlowereleven}{22}
\newcommand{\nlowertwelve}{46}
\newcommand{\nsettled}{2537}
\newcommand{\nsettledpinch}{2177}
\newcommand{\nsettledlower}{163}
\newcommand{\nsettledboth}{197}
\newcommand{\nsettledone}{520}
\newcommand{\nsettledtwo}{1995}
\newcommand{\nsettledthree}{22}
\newcommand{\nnarrowed}{376}
\newcommand{\nopen}{988}
\newcommand{\countstable}{\begin{tabular}{rlrrrr}
\toprule
crossings & KnotInfo & knots & settled & narrowed & open\\
\midrule
11 & $[1,\,2]$ & 14 & 8 & 0 & 6\\
11 & $[1,\,3]$ & 9 & 8 & 1 & 0\\
11 & $[2,\,3]$ & 36 & 33 & 0 & 3\\
12 & $[1,\,2]$ & 87 & 43 & 0 & 44\\
12 & $[1,\,3]$ & 87 & 49 & 35 & 3\\
12 & $[2,\,3]$ & 278 & 265 & 0 & 13\\
12 & $[2,\,4]$ & 3 & 3 & 0 & 0\\
13 & $[1,\,2]$ & 890 & 431 & 0 & 459\\
13 & $[1,\,3]$ & 677 & 318 & 338 & 21\\
13 & $[1,\,4]$ & 3 & 1 & 2 & 0\\
13 & $[1,\,5]$ & 1 & 0 & 0 & 1\\
13 & $[2,\,3]$ & 1437 & 1375 & 0 & 62\\
13 & $[2,\,4]$ & 3 & 3 & 0 & 0\\
\midrule
all & & 3525 & 2537 & 376 & 612\\
\bottomrule
\end{tabular}}
\newcommand{\listone}{12a\_133\\12a\_200\\12a\_212\\12a\_218\\12a\_313\\12a\_335\\12a\_373\\12a\_390\\12a\_457\\12a\_474\\12a\_611\\12a\_639\\12a\_698\\12a\_709\\12a\_750\\12a\_818\\12a\_878\\12a\_882\\12a\_895\\12a\_916\\12a\_932\\12a\_1069\\12a\_1070\\12a\_1092\\12a\_1094\\12a\_1102\\12a\_1115\\12a\_1206\\12a\_1208\\12n\_32\\12n\_70\\12n\_147\\12n\_382\\12n\_390\\12n\_444\\12n\_495\\12n\_554\\12n\_555\\12n\_597\\12n\_632\\12n\_642\\12n\_769\\12n\_784\\12n\_831\\12n\_843\\13a\_5\\13a\_6\\13a\_12\\13a\_45\\13a\_79\\13a\_129\\13a\_131\\13a\_136\\13a\_138\\13a\_156\\13a\_158\\13a\_166\\13a\_180\\13a\_197\\13a\_224\\13a\_232\\13a\_244\\13a\_247\\13a\_253\\13a\_265\\13a\_307\\13a\_318\\13a\_338\\13a\_342\\13a\_343\\13a\_353\\13a\_365\\13a\_369\\13a\_375\\13a\_381\\13a\_392\\13a\_396\\13a\_410\\13a\_476\\13a\_505\\13a\_518\\13a\_519\\13a\_530\\13a\_537\\13a\_538\\13a\_550\\13a\_591\\13a\_610\\13a\_619\\13a\_680\\13a\_687\\13a\_699\\13a\_712\\13a\_714\\13a\_721\\13a\_726\\13a\_730\\13a\_732\\13a\_734\\13a\_776\\13a\_777\\13a\_786\\13a\_815\\13a\_823\\13a\_824\\13a\_827\\13a\_845\\13a\_861\\13a\_906\\13a\_916\\13a\_925\\13a\_988\\13a\_990\\13a\_995\\13a\_1000\\13a\_1027\\13a\_1068\\13a\_1072\\13a\_1077\\13a\_1087\\13a\_1103\\13a\_1106\\13a\_1122\\13a\_1131\\13a\_1134\\13a\_1140\\13a\_1163\\13a\_1169\\13a\_1176\\13a\_1186\\13a\_1193\\13a\_1230\\13a\_1231\\13a\_1237\\13a\_1250\\13a\_1277\\13a\_1316\\13a\_1320\\13a\_1324\\13a\_1362\\13a\_1380\\13a\_1412\\13a\_1417\\13a\_1421\\13a\_1436\\13a\_1438\\13a\_1454\\13a\_1477\\13a\_1506\\13a\_1511\\13a\_1514\\13a\_1523\\13a\_1532\\13a\_1598\\13a\_1608\\13a\_1609\\13a\_1624\\13a\_1627\\13a\_1648\\13a\_1656\\13a\_1658\\13a\_1665\\13a\_1666\\13a\_1671\\13a\_1673\\13a\_1712\\13a\_1719\\13a\_1722\\13a\_1727\\13a\_1731\\13a\_1756\\13a\_1758\\13a\_1767\\13a\_1851\\13a\_1856\\13a\_1862\\13a\_1907\\13a\_2059\\13a\_2094\\13a\_2166\\13a\_2218\\13a\_2308\\13a\_2311\\13a\_2323\\13a\_2328\\13a\_2339\\13a\_2389\\13a\_2402\\13a\_2411\\13a\_2414\\13a\_2502\\13a\_2505\\13a\_2506\\13a\_2511\\13a\_2560\\13a\_2595\\13a\_2632\\13a\_2645\\13a\_2676\\13a\_2788\\13a\_2828\\13a\_2864\\13a\_2870\\13a\_2893\\13a\_2925\\13a\_3012\\13a\_3026\\13a\_3028\\13a\_3039\\13a\_3046\\13a\_3057\\13a\_3058\\13a\_3076\\13a\_3167\\13a\_3181\\13a\_3188\\13a\_3190\\13a\_3198\\13a\_3242\\13a\_3258\\13a\_3266\\13a\_3274\\13a\_3275\\13a\_3279\\13a\_3294\\13a\_3298\\13a\_3302\\13a\_3306\\13a\_3328\\13a\_3333\\13a\_3343\\13a\_3354\\13a\_3355\\13a\_3356\\13a\_3374\\13a\_3399\\13a\_3403\\13a\_3408\\13a\_3413\\13a\_3415\\13a\_3417\\13a\_3418\\13a\_3462\\13a\_3463\\13a\_3475\\13a\_3477\\13a\_3491\\13a\_3494\\13a\_3497\\13a\_3505\\13a\_3515\\13a\_3522\\13a\_3556\\13a\_3579\\13a\_3598\\13a\_3609\\13a\_3611\\13a\_3618\\13a\_3653\\13a\_3657\\13a\_3674\\13a\_3689\\13a\_3692\\13a\_3741\\13a\_3758\\13a\_3764\\13a\_3772\\13a\_3776\\13a\_3786\\13a\_3787\\13a\_3799\\13a\_3805\\13a\_3807\\13a\_3811\\13a\_3866\\13a\_3905\\13a\_3928\\13a\_3930\\13a\_3955\\13a\_3993\\13a\_4054\\13a\_4077\\13a\_4103\\13a\_4119\\13a\_4141\\13a\_4152\\13a\_4178\\13a\_4195\\13a\_4198\\13a\_4211\\13a\_4220\\13a\_4247\\13a\_4249\\13a\_4268\\13a\_4289\\13a\_4338\\13a\_4342\\13a\_4365\\13a\_4367\\13a\_4394\\13a\_4398\\13a\_4403\\13a\_4412\\13a\_4419\\13a\_4426\\13a\_4430\\13a\_4437\\13a\_4443\\13a\_4449\\13a\_4475\\13a\_4479\\13a\_4486\\13a\_4487\\13a\_4493\\13a\_4511\\13a\_4513\\13a\_4524\\13a\_4544\\13a\_4577\\13a\_4591\\13a\_4603\\13a\_4620\\13a\_4629\\13a\_4647\\13a\_4648\\13a\_4654\\13a\_4696\\13a\_4697\\13a\_4705\\13a\_4743\\13a\_4744\\13a\_4746\\13a\_4794\\13a\_4848\\13a\_4874\\13n\_7\\13n\_30\\13n\_50\\13n\_54\\13n\_78\\13n\_149\\13n\_153\\13n\_157\\13n\_231\\13n\_290\\13n\_330\\13n\_341\\13n\_344\\13n\_347\\13n\_360\\13n\_370\\13n\_406\\13n\_418\\13n\_428\\13n\_520\\13n\_529\\13n\_530\\13n\_561\\13n\_563\\13n\_620\\13n\_635\\13n\_671\\13n\_813\\13n\_816\\13n\_839\\13n\_852\\13n\_855\\13n\_892\\13n\_932\\13n\_988\\13n\_994\\13n\_999\\13n\_1005\\13n\_1057\\13n\_1061\\13n\_1062\\13n\_1067\\13n\_1111\\13n\_1114\\13n\_1201\\13n\_1243\\13n\_1275\\13n\_1277\\13n\_1302\\13n\_1328\\13n\_1419\\13n\_1490\\13n\_1492\\13n\_1494\\13n\_1500\\13n\_1582\\13n\_1585\\13n\_1609\\13n\_1628\\13n\_1638\\13n\_1670\\13n\_1677\\13n\_1678\\13n\_1688\\13n\_1693\\13n\_1726\\13n\_1789\\13n\_1932\\13n\_1935\\13n\_1976\\13n\_1992\\13n\_2021\\13n\_2022\\13n\_2080\\13n\_2088\\13n\_2115\\13n\_2150\\13n\_2193\\13n\_2245\\13n\_2265\\13n\_2383\\13n\_2404\\13n\_2406\\13n\_2409\\13n\_2453\\13n\_2503\\13n\_2610\\13n\_2629\\13n\_2690\\13n\_2694\\13n\_2706\\13n\_2741\\13n\_2765\\13n\_2773\\13n\_2812\\13n\_2834\\13n\_2837\\13n\_2851\\13n\_2857\\13n\_2926\\13n\_2942\\13n\_2949\\13n\_2966\\13n\_2977\\13n\_2985\\13n\_3108\\13n\_3190\\13n\_3198\\13n\_3236\\13n\_3311\\13n\_3334\\13n\_3338\\13n\_3344\\13n\_3351\\13n\_3402\\13n\_3452\\13n\_3501\\13n\_3526\\13n\_3602\\13n\_3641\\13n\_3667\\13n\_3695\\13n\_3723\\13n\_3739\\13n\_3740\\13n\_3745\\13n\_3752\\13n\_3771\\13n\_3776\\13n\_3899\\13n\_3902\\13n\_3916\\13n\_3919\\13n\_3921\\13n\_3923\\13n\_3933\\13n\_3968\\13n\_4023\\13n\_4030\\13n\_4107\\13n\_4115\\13n\_4152\\13n\_4155\\13n\_4162\\13n\_4193\\13n\_4202\\13n\_4247\\13n\_4272\\13n\_4300\\13n\_4305\\13n\_4306\\13n\_4313\\13n\_4333\\13n\_4334\\13n\_4351\\13n\_4354\\13n\_4380\\13n\_4457\\13n\_4477\\13n\_4485\\13n\_4509\\13n\_4515\\13n\_4517\\13n\_4521\\13n\_4543\\13n\_4614\\13n\_4632\\13n\_4642\\13n\_4657\\13n\_4666\\13n\_4709\\13n\_4746\\13n\_4870\\13n\_4896\\13n\_4940\\13n\_5006\\13n\_5009\\13n\_5012\\13n\_5015\\13n\_5041\\13n\_5067\\13n\_5089\\13n\_5095\\13n\_5104\\13n\_5108}
\newcommand{\listtwo}{11a\_6\\11a\_17\\11a\_20\\11a\_29\\11a\_49\\11a\_51\\11a\_53\\11a\_59\\11a\_64\\11a\_66\\11a\_89\\11a\_144\\11a\_161\\11a\_163\\11a\_173\\11a\_178\\11a\_188\\11a\_192\\11a\_203\\11a\_206\\11a\_214\\11a\_240\\11a\_241\\11a\_248\\11a\_259\\11a\_264\\11a\_274\\11a\_277\\11a\_278\\11a\_280\\11a\_292\\11a\_293\\11a\_297\\11a\_299\\11a\_305\\11a\_308\\11a\_314\\11a\_318\\11a\_319\\11a\_329\\11a\_337\\11a\_338\\11a\_339\\11a\_352\\11a\_354\\11a\_355\\11a\_360\\11a\_362\\11a\_366\\12a\_9\\12a\_29\\12a\_35\\12a\_36\\12a\_39\\12a\_43\\12a\_44\\12a\_47\\12a\_49\\12a\_52\\12a\_55\\12a\_64\\12a\_66\\12a\_67\\12a\_76\\12a\_85\\12a\_88\\12a\_96\\12a\_97\\12a\_103\\12a\_113\\12a\_114\\12a\_117\\12a\_134\\12a\_136\\12a\_137\\12a\_143\\12a\_146\\12a\_148\\12a\_152\\12a\_156\\12a\_167\\12a\_176\\12a\_186\\12a\_188\\12a\_193\\12a\_195\\12a\_224\\12a\_239\\12a\_254\\12a\_264\\12a\_269\\12a\_276\\12a\_277\\12a\_289\\12a\_296\\12a\_297\\12a\_304\\12a\_331\\12a\_332\\12a\_338\\12a\_347\\12a\_370\\12a\_371\\12a\_389\\12a\_392\\12a\_426\\12a\_428\\12a\_437\\12a\_446\\12a\_449\\12a\_460\\12a\_461\\12a\_480\\12a\_493\\12a\_502\\12a\_503\\12a\_507\\12a\_508\\12a\_514\\12a\_523\\12a\_533\\12a\_542\\12a\_546\\12a\_551\\12a\_553\\12a\_556\\12a\_570\\12a\_575\\12a\_587\\12a\_590\\12a\_592\\12a\_595\\12a\_597\\12a\_610\\12a\_616\\12a\_628\\12a\_633\\12a\_636\\12a\_641\\12a\_647\\12a\_650\\12a\_651\\12a\_653\\12a\_655\\12a\_657\\12a\_661\\12a\_665\\12a\_669\\12a\_674\\12a\_681\\12a\_686\\12a\_689\\12a\_692\\12a\_693\\12a\_694\\12a\_702\\12a\_703\\12a\_706\\12a\_712\\12a\_714\\12a\_726\\12a\_730\\12a\_738\\12a\_742\\12a\_753\\12a\_762\\12a\_763\\12a\_767\\12a\_783\\12a\_787\\12a\_792\\12a\_795\\12a\_796\\12a\_801\\12a\_802\\12a\_806\\12a\_811\\12a\_812\\12a\_813\\12a\_815\\12a\_816\\12a\_817\\12a\_820\\12a\_822\\12a\_827\\12a\_828\\12a\_829\\12a\_832\\12a\_840\\12a\_841\\12a\_844\\12a\_846\\12a\_853\\12a\_855\\12a\_859\\12a\_869\\12a\_873\\12a\_876\\12a\_880\\12a\_881\\12a\_885\\12a\_890\\12a\_899\\12a\_905\\12a\_912\\12a\_923\\12a\_927\\12a\_937\\12a\_946\\12a\_950\\12a\_969\\12a\_974\\12a\_981\\12a\_984\\12a\_986\\12a\_994\\12a\_996\\12a\_1001\\12a\_1002\\12a\_1004\\12a\_1006\\12a\_1009\\12a\_1012\\12a\_1013\\12a\_1015\\12a\_1021\\12a\_1025\\12a\_1030\\12a\_1031\\12a\_1035\\12a\_1036\\12a\_1037\\12a\_1044\\12a\_1049\\12a\_1052\\12a\_1058\\12a\_1059\\12a\_1079\\12a\_1080\\12a\_1084\\12a\_1088\\12a\_1091\\12a\_1097\\12a\_1099\\12a\_1103\\12a\_1106\\12a\_1107\\12a\_1112\\12a\_1114\\12a\_1116\\12a\_1120\\12a\_1122\\12a\_1137\\12a\_1142\\12a\_1143\\12a\_1152\\12a\_1159\\12a\_1164\\12a\_1171\\12a\_1172\\12a\_1177\\12a\_1179\\12a\_1180\\12a\_1184\\12a\_1189\\12a\_1193\\12a\_1197\\12a\_1200\\12a\_1222\\12a\_1223\\12a\_1229\\12a\_1232\\12a\_1240\\12a\_1242\\12a\_1244\\12a\_1248\\12a\_1252\\12a\_1255\\12a\_1256\\12a\_1260\\12a\_1265\\12a\_1266\\12a\_1270\\12a\_1274\\12a\_1276\\12a\_1279\\12a\_1288\\12n\_27\\12n\_33\\12n\_52\\12n\_67\\12n\_104\\12n\_107\\12n\_148\\12n\_153\\12n\_164\\12n\_169\\12n\_173\\12n\_174\\12n\_185\\12n\_186\\12n\_190\\12n\_211\\12n\_229\\12n\_243\\12n\_247\\12n\_260\\12n\_269\\12n\_289\\12n\_305\\12n\_317\\12n\_330\\12n\_361\\12n\_402\\12n\_432\\12n\_472\\12n\_503\\12n\_509\\12n\_517\\12n\_542\\12n\_557\\12n\_585\\12n\_603\\12n\_627\\12n\_638\\12n\_639\\12n\_651\\12n\_671\\12n\_672\\12n\_675\\12n\_682\\12n\_688\\12n\_691\\12n\_692\\12n\_694\\12n\_707\\12n\_712\\12n\_717\\12n\_722\\12n\_734\\12n\_743\\12n\_746\\12n\_755\\12n\_763\\12n\_788\\12n\_796\\12n\_803\\12n\_805\\12n\_818\\12n\_820\\12n\_827\\12n\_830\\12n\_863\\12n\_866\\12n\_874\\12n\_880\\12n\_881\\13a\_4\\13a\_8\\13a\_16\\13a\_25\\13a\_26\\13a\_29\\13a\_32\\13a\_33\\13a\_35\\13a\_40\\13a\_51\\13a\_52\\13a\_56\\13a\_58\\13a\_59\\13a\_65\\13a\_71\\13a\_82\\13a\_84\\13a\_85\\13a\_88\\13a\_89\\13a\_96\\13a\_97\\13a\_98\\13a\_103\\13a\_112\\13a\_122\\13a\_123\\13a\_125\\13a\_132\\13a\_133\\13a\_134\\13a\_140\\13a\_150\\13a\_151\\13a\_154\\13a\_157\\13a\_160\\13a\_162\\13a\_170\\13a\_171\\13a\_173\\13a\_185\\13a\_189\\13a\_191\\13a\_193\\13a\_195\\13a\_201\\13a\_203\\13a\_209\\13a\_214\\13a\_217\\13a\_223\\13a\_225\\13a\_231\\13a\_236\\13a\_251\\13a\_254\\13a\_257\\13a\_273\\13a\_275\\13a\_279\\13a\_284\\13a\_286\\13a\_295\\13a\_304\\13a\_309\\13a\_314\\13a\_316\\13a\_319\\13a\_333\\13a\_334\\13a\_345\\13a\_349\\13a\_355\\13a\_361\\13a\_362\\13a\_366\\13a\_370\\13a\_374\\13a\_382\\13a\_393\\13a\_395\\13a\_402\\13a\_403\\13a\_406\\13a\_408\\13a\_416\\13a\_423\\13a\_426\\13a\_428\\13a\_429\\13a\_433\\13a\_435\\13a\_437\\13a\_438\\13a\_443\\13a\_445\\13a\_448\\13a\_451\\13a\_468\\13a\_469\\13a\_473\\13a\_474\\13a\_480\\13a\_481\\13a\_496\\13a\_497\\13a\_501\\13a\_508\\13a\_513\\13a\_525\\13a\_535\\13a\_536\\13a\_539\\13a\_540\\13a\_542\\13a\_544\\13a\_545\\13a\_546\\13a\_551\\13a\_553\\13a\_561\\13a\_564\\13a\_569\\13a\_571\\13a\_573\\13a\_574\\13a\_584\\13a\_589\\13a\_590\\13a\_593\\13a\_594\\13a\_598\\13a\_599\\13a\_603\\13a\_606\\13a\_609\\13a\_616\\13a\_618\\13a\_620\\13a\_621\\13a\_622\\13a\_623\\13a\_625\\13a\_627\\13a\_629\\13a\_631\\13a\_635\\13a\_636\\13a\_643\\13a\_644\\13a\_648\\13a\_652\\13a\_654\\13a\_656\\13a\_658\\13a\_661\\13a\_662\\13a\_663\\13a\_665\\13a\_667\\13a\_669\\13a\_670\\13a\_675\\13a\_676\\13a\_679\\13a\_681\\13a\_682\\13a\_686\\13a\_702\\13a\_705\\13a\_716\\13a\_718\\13a\_720\\13a\_738\\13a\_740\\13a\_744\\13a\_754\\13a\_762\\13a\_765\\13a\_768\\13a\_774\\13a\_784\\13a\_789\\13a\_808\\13a\_813\\13a\_819\\13a\_826\\13a\_831\\13a\_832\\13a\_855\\13a\_866\\13a\_867\\13a\_892\\13a\_895\\13a\_896\\13a\_900\\13a\_905\\13a\_910\\13a\_918\\13a\_924\\13a\_931\\13a\_933\\13a\_938\\13a\_945\\13a\_955\\13a\_964\\13a\_966\\13a\_968\\13a\_971\\13a\_976\\13a\_980\\13a\_989\\13a\_991\\13a\_993\\13a\_1002\\13a\_1020\\13a\_1022\\13a\_1041\\13a\_1043\\13a\_1045\\13a\_1050\\13a\_1051\\13a\_1063\\13a\_1066\\13a\_1069\\13a\_1070\\13a\_1089\\13a\_1099\\13a\_1100\\13a\_1102\\13a\_1120\\13a\_1129\\13a\_1133\\13a\_1137\\13a\_1142\\13a\_1147\\13a\_1149\\13a\_1157\\13a\_1158\\13a\_1182\\13a\_1183\\13a\_1188\\13a\_1192\\13a\_1196\\13a\_1203\\13a\_1221\\13a\_1226\\13a\_1233\\13a\_1235\\13a\_1243\\13a\_1245\\13a\_1246\\13a\_1255\\13a\_1258\\13a\_1259\\13a\_1266\\13a\_1283\\13a\_1302\\13a\_1306\\13a\_1331\\13a\_1332\\13a\_1335\\13a\_1349\\13a\_1353\\13a\_1355\\13a\_1359\\13a\_1371\\13a\_1372\\13a\_1377\\13a\_1378\\13a\_1381\\13a\_1395\\13a\_1396\\13a\_1400\\13a\_1404\\13a\_1405\\13a\_1406\\13a\_1414\\13a\_1419\\13a\_1429\\13a\_1440\\13a\_1442\\13a\_1452\\13a\_1461\\13a\_1462\\13a\_1464\\13a\_1465\\13a\_1466\\13a\_1467\\13a\_1468\\13a\_1472\\13a\_1474\\13a\_1475\\13a\_1482\\13a\_1488\\13a\_1492\\13a\_1495\\13a\_1503\\13a\_1520\\13a\_1524\\13a\_1541\\13a\_1542\\13a\_1543\\13a\_1544\\13a\_1546\\13a\_1547\\13a\_1550\\13a\_1553\\13a\_1566\\13a\_1567\\13a\_1573\\13a\_1580\\13a\_1583\\13a\_1590\\13a\_1601\\13a\_1603\\13a\_1614\\13a\_1617\\13a\_1618\\13a\_1621\\13a\_1629\\13a\_1630\\13a\_1634\\13a\_1640\\13a\_1641\\13a\_1643\\13a\_1645\\13a\_1652\\13a\_1660\\13a\_1682\\13a\_1683\\13a\_1692\\13a\_1697\\13a\_1698\\13a\_1703\\13a\_1706\\13a\_1718\\13a\_1720\\13a\_1721\\13a\_1725\\13a\_1740\\13a\_1751\\13a\_1757\\13a\_1762\\13a\_1785\\13a\_1786\\13a\_1791\\13a\_1794\\13a\_1820\\13a\_1825\\13a\_1827\\13a\_1829\\13a\_1854\\13a\_1855\\13a\_1859\\13a\_1864\\13a\_1865\\13a\_1867\\13a\_1879\\13a\_1885\\13a\_1895\\13a\_1897\\13a\_1899\\13a\_1901\\13a\_1922\\13a\_1927\\13a\_1928\\13a\_1929\\13a\_1943\\13a\_1946\\13a\_1952\\13a\_1955\\13a\_1957\\13a\_1960\\13a\_1961\\13a\_1969\\13a\_1971\\13a\_1972\\13a\_1976\\13a\_1982\\13a\_1984\\13a\_1985\\13a\_1987\\13a\_1989\\13a\_1990\\13a\_1997\\13a\_1999\\13a\_2008\\13a\_2018\\13a\_2020\\13a\_2025\\13a\_2030\\13a\_2032\\13a\_2037\\13a\_2038\\13a\_2042\\13a\_2052\\13a\_2054\\13a\_2055\\13a\_2064\\13a\_2076\\13a\_2077\\13a\_2078\\13a\_2081\\13a\_2088\\13a\_2092\\13a\_2093\\13a\_2097\\13a\_2098\\13a\_2107\\13a\_2109\\13a\_2114\\13a\_2119\\13a\_2124\\13a\_2129\\13a\_2132\\13a\_2135\\13a\_2139\\13a\_2147\\13a\_2154\\13a\_2157\\13a\_2159\\13a\_2167\\13a\_2173\\13a\_2175\\13a\_2185\\13a\_2194\\13a\_2202\\13a\_2207\\13a\_2228\\13a\_2232\\13a\_2238\\13a\_2245\\13a\_2249\\13a\_2255\\13a\_2263\\13a\_2265\\13a\_2268\\13a\_2269\\13a\_2270\\13a\_2271\\13a\_2284\\13a\_2286\\13a\_2287\\13a\_2290\\13a\_2307\\13a\_2310\\13a\_2315\\13a\_2319\\13a\_2326\\13a\_2333\\13a\_2335\\13a\_2359\\13a\_2360\\13a\_2362\\13a\_2367\\13a\_2370\\13a\_2372\\13a\_2375\\13a\_2376\\13a\_2381\\13a\_2386\\13a\_2387\\13a\_2390\\13a\_2395\\13a\_2399\\13a\_2400\\13a\_2405\\13a\_2407\\13a\_2408\\13a\_2409\\13a\_2417\\13a\_2419\\13a\_2421\\13a\_2428\\13a\_2432\\13a\_2434\\13a\_2438\\13a\_2443\\13a\_2445\\13a\_2450\\13a\_2452\\13a\_2453\\13a\_2461\\13a\_2462\\13a\_2474\\13a\_2476\\13a\_2478\\13a\_2484\\13a\_2485\\13a\_2497\\13a\_2501\\13a\_2520\\13a\_2526\\13a\_2533\\13a\_2536\\13a\_2538\\13a\_2542\\13a\_2543\\13a\_2544\\13a\_2548\\13a\_2550\\13a\_2552\\13a\_2554\\13a\_2561\\13a\_2564\\13a\_2565\\13a\_2578\\13a\_2593\\13a\_2598\\13a\_2600\\13a\_2604\\13a\_2606\\13a\_2607\\13a\_2610\\13a\_2611\\13a\_2616\\13a\_2620\\13a\_2628\\13a\_2636\\13a\_2641\\13a\_2642\\13a\_2644\\13a\_2647\\13a\_2655\\13a\_2656\\13a\_2667\\13a\_2672\\13a\_2674\\13a\_2679\\13a\_2683\\13a\_2690\\13a\_2697\\13a\_2706\\13a\_2709\\13a\_2711\\13a\_2717\\13a\_2721\\13a\_2722\\13a\_2724\\13a\_2732\\13a\_2735\\13a\_2736\\13a\_2742\\13a\_2751\\13a\_2757\\13a\_2761\\13a\_2763\\13a\_2767\\13a\_2769\\13a\_2772\\13a\_2782\\13a\_2784\\13a\_2786\\13a\_2790\\13a\_2810\\13a\_2813\\13a\_2814\\13a\_2816\\13a\_2821\\13a\_2823\\13a\_2824\\13a\_2829\\13a\_2830\\13a\_2833\\13a\_2835\\13a\_2840\\13a\_2843\\13a\_2844\\13a\_2853\\13a\_2856\\13a\_2863\\13a\_2865\\13a\_2866\\13a\_2872\\13a\_2879\\13a\_2880\\13a\_2881\\13a\_2884\\13a\_2886\\13a\_2888\\13a\_2897\\13a\_2911\\13a\_2920\\13a\_2933\\13a\_2937\\13a\_2943\\13a\_2944\\13a\_2945\\13a\_2947\\13a\_2959\\13a\_2970\\13a\_2981\\13a\_2982\\13a\_2984\\13a\_2994\\13a\_2996\\13a\_3000\\13a\_3007\\13a\_3016\\13a\_3018\\13a\_3021\\13a\_3022\\13a\_3030\\13a\_3036\\13a\_3040\\13a\_3042\\13a\_3044\\13a\_3045\\13a\_3059\\13a\_3060\\13a\_3062\\13a\_3068\\13a\_3074\\13a\_3080\\13a\_3083\\13a\_3089\\13a\_3090\\13a\_3093\\13a\_3096\\13a\_3099\\13a\_3100\\13a\_3110\\13a\_3112\\13a\_3114\\13a\_3116\\13a\_3120\\13a\_3122\\13a\_3127\\13a\_3128\\13a\_3129\\13a\_3130\\13a\_3131\\13a\_3132\\13a\_3133\\13a\_3139\\13a\_3141\\13a\_3143\\13a\_3146\\13a\_3153\\13a\_3155\\13a\_3156\\13a\_3159\\13a\_3160\\13a\_3161\\13a\_3164\\13a\_3166\\13a\_3174\\13a\_3176\\13a\_3178\\13a\_3179\\13a\_3182\\13a\_3185\\13a\_3192\\13a\_3207\\13a\_3212\\13a\_3217\\13a\_3222\\13a\_3223\\13a\_3225\\13a\_3230\\13a\_3232\\13a\_3234\\13a\_3235\\13a\_3237\\13a\_3239\\13a\_3240\\13a\_3244\\13a\_3248\\13a\_3250\\13a\_3253\\13a\_3255\\13a\_3257\\13a\_3259\\13a\_3261\\13a\_3262\\13a\_3265\\13a\_3267\\13a\_3277\\13a\_3282\\13a\_3283\\13a\_3290\\13a\_3297\\13a\_3300\\13a\_3308\\13a\_3309\\13a\_3313\\13a\_3315\\13a\_3324\\13a\_3327\\13a\_3329\\13a\_3332\\13a\_3334\\13a\_3336\\13a\_3337\\13a\_3340\\13a\_3341\\13a\_3344\\13a\_3358\\13a\_3364\\13a\_3371\\13a\_3373\\13a\_3377\\13a\_3379\\13a\_3380\\13a\_3386\\13a\_3387\\13a\_3388\\13a\_3390\\13a\_3400\\13a\_3401\\13a\_3426\\13a\_3428\\13a\_3436\\13a\_3444\\13a\_3454\\13a\_3459\\13a\_3460\\13a\_3461\\13a\_3468\\13a\_3474\\13a\_3478\\13a\_3479\\13a\_3482\\13a\_3492\\13a\_3500\\13a\_3504\\13a\_3508\\13a\_3512\\13a\_3524\\13a\_3530\\13a\_3531\\13a\_3534\\13a\_3535\\13a\_3537\\13a\_3543\\13a\_3551\\13a\_3553\\13a\_3554\\13a\_3558\\13a\_3559\\13a\_3564\\13a\_3567\\13a\_3570\\13a\_3577\\13a\_3589\\13a\_3592\\13a\_3594\\13a\_3595\\13a\_3596\\13a\_3597\\13a\_3605\\13a\_3608\\13a\_3615\\13a\_3621\\13a\_3631\\13a\_3639\\13a\_3641\\13a\_3642\\13a\_3652\\13a\_3660\\13a\_3663\\13a\_3667\\13a\_3669\\13a\_3670\\13a\_3676\\13a\_3686\\13a\_3687\\13a\_3688\\13a\_3690\\13a\_3691\\13a\_3699\\13a\_3704\\13a\_3706\\13a\_3714\\13a\_3716\\13a\_3718\\13a\_3720\\13a\_3731\\13a\_3736\\13a\_3737\\13a\_3743\\13a\_3745\\13a\_3749\\13a\_3752\\13a\_3754\\13a\_3759\\13a\_3792\\13a\_3794\\13a\_3795\\13a\_3796\\13a\_3800\\13a\_3804\\13a\_3812\\13a\_3822\\13a\_3823\\13a\_3827\\13a\_3841\\13a\_3844\\13a\_3845\\13a\_3847\\13a\_3852\\13a\_3858\\13a\_3864\\13a\_3865\\13a\_3869\\13a\_3872\\13a\_3880\\13a\_3883\\13a\_3888\\13a\_3892\\13a\_3896\\13a\_3899\\13a\_3901\\13a\_3912\\13a\_3913\\13a\_3917\\13a\_3922\\13a\_3926\\13a\_3927\\13a\_3931\\13a\_3940\\13a\_3943\\13a\_3956\\13a\_3960\\13a\_3961\\13a\_3962\\13a\_3972\\13a\_3975\\13a\_3976\\13a\_3977\\13a\_3978\\13a\_3979\\13a\_3990\\13a\_3998\\13a\_3999\\13a\_4005\\13a\_4009\\13a\_4012\\13a\_4017\\13a\_4021\\13a\_4025\\13a\_4034\\13a\_4039\\13a\_4041\\13a\_4043\\13a\_4046\\13a\_4047\\13a\_4048\\13a\_4049\\13a\_4050\\13a\_4053\\13a\_4055\\13a\_4057\\13a\_4059\\13a\_4065\\13a\_4069\\13a\_4072\\13a\_4080\\13a\_4083\\13a\_4087\\13a\_4088\\13a\_4089\\13a\_4090\\13a\_4093\\13a\_4096\\13a\_4101\\13a\_4106\\13a\_4108\\13a\_4109\\13a\_4112\\13a\_4113\\13a\_4116\\13a\_4118\\13a\_4120\\13a\_4123\\13a\_4125\\13a\_4127\\13a\_4128\\13a\_4129\\13a\_4134\\13a\_4149\\13a\_4153\\13a\_4155\\13a\_4160\\13a\_4162\\13a\_4163\\13a\_4170\\13a\_4176\\13a\_4179\\13a\_4182\\13a\_4184\\13a\_4190\\13a\_4191\\13a\_4193\\13a\_4201\\13a\_4206\\13a\_4207\\13a\_4208\\13a\_4210\\13a\_4214\\13a\_4216\\13a\_4225\\13a\_4226\\13a\_4229\\13a\_4231\\13a\_4233\\13a\_4235\\13a\_4241\\13a\_4246\\13a\_4250\\13a\_4254\\13a\_4255\\13a\_4258\\13a\_4259\\13a\_4260\\13a\_4263\\13a\_4265\\13a\_4281\\13a\_4286\\13a\_4287\\13a\_4290\\13a\_4293\\13a\_4304\\13a\_4306\\13a\_4307\\13a\_4309\\13a\_4310\\13a\_4312\\13a\_4313\\13a\_4316\\13a\_4320\\13a\_4321\\13a\_4322\\13a\_4328\\13a\_4329\\13a\_4331\\13a\_4345\\13a\_4352\\13a\_4353\\13a\_4359\\13a\_4362\\13a\_4368\\13a\_4369\\13a\_4376\\13a\_4377\\13a\_4379\\13a\_4381\\13a\_4382\\13a\_4385\\13a\_4386\\13a\_4388\\13a\_4390\\13a\_4393\\13a\_4409\\13a\_4421\\13a\_4422\\13a\_4424\\13a\_4427\\13a\_4429\\13a\_4431\\13a\_4445\\13a\_4454\\13a\_4457\\13a\_4459\\13a\_4464\\13a\_4468\\13a\_4478\\13a\_4480\\13a\_4481\\13a\_4482\\13a\_4483\\13a\_4484\\13a\_4488\\13a\_4492\\13a\_4499\\13a\_4502\\13a\_4503\\13a\_4504\\13a\_4508\\13a\_4509\\13a\_4517\\13a\_4518\\13a\_4531\\13a\_4535\\13a\_4536\\13a\_4538\\13a\_4541\\13a\_4549\\13a\_4551\\13a\_4552\\13a\_4553\\13a\_4554\\13a\_4555\\13a\_4557\\13a\_4558\\13a\_4565\\13a\_4567\\13a\_4568\\13a\_4571\\13a\_4574\\13a\_4576\\13a\_4578\\13a\_4580\\13a\_4589\\13a\_4595\\13a\_4599\\13a\_4606\\13a\_4607\\13a\_4609\\13a\_4616\\13a\_4622\\13a\_4624\\13a\_4631\\13a\_4636\\13a\_4643\\13a\_4644\\13a\_4645\\13a\_4646\\13a\_4656\\13a\_4657\\13a\_4658\\13a\_4660\\13a\_4667\\13a\_4673\\13a\_4691\\13a\_4695\\13a\_4698\\13a\_4704\\13a\_4706\\13a\_4707\\13a\_4709\\13a\_4718\\13a\_4725\\13a\_4732\\13a\_4733\\13a\_4737\\13a\_4739\\13a\_4740\\13a\_4745\\13a\_4747\\13a\_4757\\13a\_4758\\13a\_4764\\13a\_4766\\13a\_4767\\13a\_4768\\13a\_4770\\13a\_4771\\13a\_4774\\13a\_4787\\13a\_4790\\13a\_4792\\13a\_4796\\13a\_4798\\13a\_4807\\13a\_4809\\13a\_4811\\13a\_4812\\13a\_4813\\13a\_4817\\13a\_4820\\13a\_4823\\13a\_4824\\13a\_4826\\13a\_4828\\13a\_4832\\13a\_4835\\13a\_4837\\13a\_4840\\13a\_4842\\13a\_4843\\13a\_4844\\13a\_4849\\13a\_4850\\13a\_4851\\13a\_4853\\13a\_4857\\13a\_4858\\13a\_4859\\13a\_4861\\13a\_4862\\13a\_4864\\13a\_4865\\13a\_4869\\13a\_4873\\13a\_4875\\13a\_4876\\13n\_6\\13n\_17\\13n\_25\\13n\_32\\13n\_41\\13n\_44\\13n\_45\\13n\_75\\13n\_82\\13n\_113\\13n\_131\\13n\_150\\13n\_154\\13n\_163\\13n\_177\\13n\_186\\13n\_187\\13n\_204\\13n\_215\\13n\_222\\13n\_225\\13n\_241\\13n\_247\\13n\_265\\13n\_274\\13n\_285\\13n\_300\\13n\_312\\13n\_317\\13n\_321\\13n\_338\\13n\_339\\13n\_342\\13n\_356\\13n\_368\\13n\_378\\13n\_384\\13n\_393\\13n\_429\\13n\_437\\13n\_438\\13n\_458\\13n\_486\\13n\_495\\13n\_500\\13n\_503\\13n\_517\\13n\_519\\13n\_535\\13n\_546\\13n\_550\\13n\_569\\13n\_596\\13n\_601\\13n\_616\\13n\_618\\13n\_622\\13n\_629\\13n\_631\\13n\_645\\13n\_663\\13n\_670\\13n\_675\\13n\_679\\13n\_680\\13n\_681\\13n\_691\\13n\_709\\13n\_719\\13n\_723\\13n\_732\\13n\_744\\13n\_747\\13n\_755\\13n\_767\\13n\_771\\13n\_773\\13n\_776\\13n\_787\\13n\_832\\13n\_841\\13n\_858\\13n\_895\\13n\_900\\13n\_907\\13n\_914\\13n\_916\\13n\_929\\13n\_935\\13n\_943\\13n\_947\\13n\_961\\13n\_962\\13n\_965\\13n\_974\\13n\_981\\13n\_983\\13n\_1004\\13n\_1008\\13n\_1011\\13n\_1029\\13n\_1043\\13n\_1049\\13n\_1058\\13n\_1063\\13n\_1072\\13n\_1075\\13n\_1077\\13n\_1080\\13n\_1086\\13n\_1095\\13n\_1097\\13n\_1104\\13n\_1106\\13n\_1115\\13n\_1117\\13n\_1176\\13n\_1180\\13n\_1187\\13n\_1192\\13n\_1204\\13n\_1205\\13n\_1212\\13n\_1217\\13n\_1222\\13n\_1225\\13n\_1229\\13n\_1235\\13n\_1238\\13n\_1239\\13n\_1240\\13n\_1273\\13n\_1278\\13n\_1294\\13n\_1300\\13n\_1323\\13n\_1329\\13n\_1332\\13n\_1335\\13n\_1341\\13n\_1344\\13n\_1382\\13n\_1383\\13n\_1390\\13n\_1406\\13n\_1415\\13n\_1417\\13n\_1421\\13n\_1439\\13n\_1498\\13n\_1503\\13n\_1504\\13n\_1508\\13n\_1516\\13n\_1517\\13n\_1519\\13n\_1520\\13n\_1528\\13n\_1533\\13n\_1535\\13n\_1567\\13n\_1590\\13n\_1591\\13n\_1592\\13n\_1594\\13n\_1595\\13n\_1643\\13n\_1654\\13n\_1656\\13n\_1660\\13n\_1679\\13n\_1681\\13n\_1684\\13n\_1703\\13n\_1730\\13n\_1732\\13n\_1771\\13n\_1784\\13n\_1791\\13n\_1838\\13n\_1851\\13n\_1858\\13n\_1892\\13n\_1906\\13n\_1907\\13n\_1927\\13n\_1931\\13n\_1933\\13n\_1936\\13n\_1950\\13n\_1998\\13n\_2042\\13n\_2061\\13n\_2065\\13n\_2089\\13n\_2104\\13n\_2110\\13n\_2114\\13n\_2121\\13n\_2123\\13n\_2125\\13n\_2131\\13n\_2143\\13n\_2159\\13n\_2171\\13n\_2190\\13n\_2192\\13n\_2195\\13n\_2231\\13n\_2248\\13n\_2252\\13n\_2253\\13n\_2280\\13n\_2287\\13n\_2297\\13n\_2314\\13n\_2335\\13n\_2345\\13n\_2346\\13n\_2362\\13n\_2368\\13n\_2372\\13n\_2421\\13n\_2426\\13n\_2437\\13n\_2442\\13n\_2487\\13n\_2515\\13n\_2530\\13n\_2536\\13n\_2561\\13n\_2572\\13n\_2586\\13n\_2603\\13n\_2608\\13n\_2649\\13n\_2657\\13n\_2661\\13n\_2664\\13n\_2672\\13n\_2677\\13n\_2689\\13n\_2711\\13n\_2719\\13n\_2722\\13n\_2751\\13n\_2775\\13n\_2805\\13n\_2816\\13n\_2827\\13n\_2870\\13n\_2881\\13n\_2883\\13n\_2898\\13n\_2899\\13n\_2900\\13n\_2901\\13n\_2907\\13n\_2920\\13n\_2929\\13n\_2935\\13n\_2939\\13n\_2959\\13n\_2972\\13n\_2981\\13n\_2995\\13n\_3000\\13n\_3003\\13n\_3010\\13n\_3015\\13n\_3034\\13n\_3038\\13n\_3046\\13n\_3048\\13n\_3051\\13n\_3060\\13n\_3063\\13n\_3068\\13n\_3083\\13n\_3103\\13n\_3111\\13n\_3113\\13n\_3124\\13n\_3137\\13n\_3145\\13n\_3156\\13n\_3188\\13n\_3199\\13n\_3200\\13n\_3210\\13n\_3212\\13n\_3216\\13n\_3239\\13n\_3248\\13n\_3254\\13n\_3258\\13n\_3274\\13n\_3276\\13n\_3293\\13n\_3294\\13n\_3305\\13n\_3333\\13n\_3337\\13n\_3354\\13n\_3368\\13n\_3419\\13n\_3426\\13n\_3427\\13n\_3432\\13n\_3435\\13n\_3447\\13n\_3454\\13n\_3461\\13n\_3470\\13n\_3476\\13n\_3484\\13n\_3486\\13n\_3491\\13n\_3493\\13n\_3497\\13n\_3499\\13n\_3502\\13n\_3508\\13n\_3515\\13n\_3518\\13n\_3520\\13n\_3532\\13n\_3566\\13n\_3576\\13n\_3579\\13n\_3583\\13n\_3589\\13n\_3593\\13n\_3595\\13n\_3608\\13n\_3633\\13n\_3643\\13n\_3648\\13n\_3653\\13n\_3657\\13n\_3673\\13n\_3676\\13n\_3697\\13n\_3698\\13n\_3709\\13n\_3712\\13n\_3714\\13n\_3721\\13n\_3724\\13n\_3729\\13n\_3735\\13n\_3736\\13n\_3743\\13n\_3749\\13n\_3774\\13n\_3780\\13n\_3785\\13n\_3788\\13n\_3799\\13n\_3801\\13n\_3806\\13n\_3807\\13n\_3810\\13n\_3813\\13n\_3822\\13n\_3824\\13n\_3836\\13n\_3837\\13n\_3849\\13n\_3852\\13n\_3854\\13n\_3873\\13n\_3879\\13n\_3882\\13n\_3907\\13n\_3913\\13n\_3920\\13n\_3928\\13n\_3942\\13n\_3990\\13n\_4008\\13n\_4015\\13n\_4029\\13n\_4043\\13n\_4050\\13n\_4054\\13n\_4057\\13n\_4073\\13n\_4085\\13n\_4094\\13n\_4095\\13n\_4129\\13n\_4133\\13n\_4137\\13n\_4164\\13n\_4165\\13n\_4169\\13n\_4179\\13n\_4181\\13n\_4183\\13n\_4186\\13n\_4187\\13n\_4189\\13n\_4200\\13n\_4207\\13n\_4208\\13n\_4218\\13n\_4221\\13n\_4223\\13n\_4243\\13n\_4256\\13n\_4261\\13n\_4263\\13n\_4269\\13n\_4278\\13n\_4289\\13n\_4314\\13n\_4341\\13n\_4342\\13n\_4343\\13n\_4353\\13n\_4359\\13n\_4360\\13n\_4361\\13n\_4362\\13n\_4363\\13n\_4365\\13n\_4366\\13n\_4373\\13n\_4377\\13n\_4402\\13n\_4404\\13n\_4408\\13n\_4410\\13n\_4418\\13n\_4447\\13n\_4455\\13n\_4463\\13n\_4469\\13n\_4501\\13n\_4520\\13n\_4537\\13n\_4546\\13n\_4550\\13n\_4556\\13n\_4561\\13n\_4565\\13n\_4567\\13n\_4568\\13n\_4569\\13n\_4576\\13n\_4588\\13n\_4589\\13n\_4599\\13n\_4605\\13n\_4615\\13n\_4622\\13n\_4633\\13n\_4637\\13n\_4641\\13n\_4644\\13n\_4649\\13n\_4650\\13n\_4656\\13n\_4660\\13n\_4673\\13n\_4676\\13n\_4679\\13n\_4687\\13n\_4693\\13n\_4700\\13n\_4704\\13n\_4714\\13n\_4717\\13n\_4729\\13n\_4730\\13n\_4731\\13n\_4739\\13n\_4743\\13n\_4747\\13n\_4751\\13n\_4752\\13n\_4761\\13n\_4767\\13n\_4774\\13n\_4782\\13n\_4791\\13n\_4793\\13n\_4795\\13n\_4804\\13n\_4813\\13n\_4820\\13n\_4822\\13n\_4837\\13n\_4846\\13n\_4847\\13n\_4853\\13n\_4865\\13n\_4867\\13n\_4868\\13n\_4879\\13n\_4884\\13n\_4887\\13n\_4890\\13n\_4891\\13n\_4899\\13n\_4900\\13n\_4906\\13n\_4911\\13n\_4914\\13n\_4928\\13n\_4931\\13n\_4933\\13n\_4936\\13n\_4941\\13n\_4943\\13n\_4949\\13n\_4951\\13n\_4952\\13n\_4960\\13n\_4966\\13n\_4967\\13n\_4972\\13n\_4975\\13n\_4978\\13n\_4981\\13n\_4982\\13n\_4985\\13n\_4990\\13n\_4991\\13n\_4993\\13n\_4994\\13n\_4997\\13n\_5003\\13n\_5013\\13n\_5014\\13n\_5021\\13n\_5022\\13n\_5030\\13n\_5044\\13n\_5046\\13n\_5054\\13n\_5070\\13n\_5075\\13n\_5084\\13n\_5100\\13n\_5102}
\newcommand{\listthree}{12a\_561\\12a\_569\\12a\_664\\12a\_780\\12a\_906\\12a\_907\\12a\_1123\\13a\_1398\\13a\_1557\\13a\_1662\\13a\_1826\\13a\_1974\\13a\_1992\\13a\_2890\\13a\_3034\\13a\_3121\\13a\_3585\\13a\_3747\\13a\_3936\\13a\_4566\\13a\_4637\\13a\_4674}
\newcommand{\selftesttable}{\begin{tabular}{lr}
\toprule
corruption & rejected\\
\midrule
other smoothing of the crossing & 20 of 20\\
next crossing of the diagram & 13 of 20\\
wrong target ($12a_{1}$) & 20 of 20\\
bound lowered by one & 20 of 20\\
\bottomrule
\end{tabular}
}


\newtheorem{theorem}{Theorem}[section]
\newtheorem{lemma}[theorem]{Lemma}
\newtheorem{corollary}[theorem]{Corollary}
\theoremstyle{definition}
\newtheorem{definition}[theorem]{Definition}
\theoremstyle{remark}
\newtheorem{remark}[theorem]{Remark}
\numberwithin{equation}{section}

\newcommand{\F}{\mathbb{F}}
\DeclareMathOperator{\rk}{rank}
\DeclareMathOperator{\Sp}{Sp}
\DeclareMathOperator{\GL}{GL}
\renewcommand{\O}{\mathrm{O}}

\title[Bridge index and nonorientable four-genus through thirteen crossings]{The bridge index and the nonorientable four-genus of knots with at most thirteen crossings}
\author{Deep Bhattacharjee}
\address{Formerly, Electro-Gravitational Space Propulsion Laboratory (EGSPL), Bhubaneswar, Odisha 751030, India; formerly, University of Calcutta, Kolkata 700073, India}
\email{itsdeep@live.com}
\date{October 22, 2026}
\subjclass[2020]{Primary 57K10; Secondary 57K40, 20C33, 11E12}
\keywords{bridge index, meridional rank, transvections, nonorientable four-genus, pinch moves, Goeritz lattices, Donaldson's theorem, KnotInfo}

\begin{document}
\begin{abstract}
We determine two invariants for knots that the KnotInfo database leaves open. KnotInfo lists
the bridge index of \btotal{} prime knots with thirteen crossings as either $3$ or $4$; we
prove that \bclosed{} of them have bridge index and meridional rank $4$, so the meridional
rank conjecture holds for them. The lower bound comes from representations of the knot
group in $\GL_4(\F_q)$ that send a meridian to a transvection or a reflection and fix no
nonzero vector. KnotInfo leaves the smooth nonorientable four-genus open for \ntotal{}
prime knots with eleven to thirteen crossings; we determine it for \nsettled{} of them and
narrow the interval for \nnarrowed{} more. Upper bounds come from pinch moves, lower bounds
from the Goeritz-lattice obstructions of Jabuka and Kelly evaluated by an exhaustive search
for lattice embeddings. Every bound is a certificate checked by a program independent of
the search that found it.
\end{abstract}
\maketitle

\section{Introduction}

The KnotInfo database \cite{KnotInfo} records, for each prime knot with at most thirteen
crossings, the values of several dozen invariants, and an interval where the value is not
known. This paper closes many of these intervals for two invariants, the bridge index and
the smooth nonorientable four-genus, both in the version 2026.10.5 of the database. The two
parts are independent and use different methods; they share the source of the diagrams,
the form of the results, a list of certificates, one per knot, each checked by a short
program, and the treatment of the knots that remain open.

\subsection*{The bridge index}
Let $K\subset S^3$ be a knot, $G=\pi_1(S^3\setminus K)$ its group and $m\in G$ a meridian.
The \emph{bridge index} $b(K)$ is the least number of local maxima of a height function on
a knot isotopic to $K$. The \emph{meridional rank} $\mu(K)$ is the least number of
meridians, that is conjugates of $m^{\pm1}$, that generate $G$. The overpasses of a
$b$-bridge presentation give $b$ meridians generating $G$, so $\mu(K)\le b(K)$. Cappell and
Shaneson asked whether equality always holds \cite[Problem~1.11]{Kirby}; this is the
meridional rank conjecture.

Blair, Kjuchukova, Velazquez and Villanueva defined the Wirtinger number $\omega(D)$ of a
diagram $D$ and proved that its minimum over the diagrams of $K$ is $b(K)$ \cite{BKVV}. In
particular
\begin{equation}\label{eq:upper}
b(K)\le\omega(D)\quad\text{for every diagram $D$ of $K$.}
\end{equation}
Blair, Kjuchukova and Morrison bounded $\mu(K)$ from below by surjections of $G$ onto
Coxeter groups sending meridians to reflections, and so computed $b(K)=\mu(K)$ for
595{,}515 knots up to sixteen crossings \cite{BKM}. KnotInfo records the bridge index of
thirteen-crossing knots from these bounds and from the plain sphere number \cite{BKP}. In
the version 2026.10.5, \btotal{} knots with thirteen crossings have the entry $[3,4]$.

\begin{theorem}\label{thm:bridge}
For each of the \bclosed{} knots $K$ listed in Appendix~\ref{app:bclosed},
\[
b(K)=\mu(K)=4 .
\]
In particular the meridional rank conjecture holds for these knots.
\end{theorem}

The upper bound is \eqref{eq:upper} applied to the KnotInfo diagram of $K$, whose Wirtinger
number is $4$. The lower bound is Lemma~\ref{lem:fixed}, applied to a representation
$\rho\colon G\to\GL_4(\F_q)$ with $q\in\{3,5\}$ that sends the meridian to a transvection of
$\Sp_4(\F_q)$ or to a reflection of an orthogonal group $\O_4(\F_5)$. Reflections in
Coxeter groups are pseudo-reflections too; the transvections, of order $3$ and $5$, are
what settles most of these knots. Table~\ref{tab:groups} gives the counts.

\begin{table}[ht]
\centering
\begin{tabular}{llr}
\toprule
group & image of a meridian & knots \\
\midrule
\bgrouptable
\bottomrule
\end{tabular}
\caption{The certificates of Theorem~\ref{thm:bridge}, by the group used.}\label{tab:groups}
\end{table}

\subsection*{The nonorientable four-genus}
Every knot $K\subset S^3=\partial B^4$ bounds a smooth surface in $B^4$; if $K$ is not
slice, it bounds nonorientable ones as well. The \emph{nonorientable four-genus}
$\gamma_4(K)$ is the least first Betti number of a smooth compact nonorientable surface
$F\subset B^4$ with $\partial F=K$; for slice knots one sets $\gamma_4(K)=1$. It was studied
by Murakami and Yasuhara \cite{MY} and by Gilmer and Livingston \cite{GiL}. Its values on
small knots were computed by Jabuka and Kelly for eight and nine crossings \cite{JK}, by
Ghanbarian for ten \cite{Gh} and by Fairchild for nonalternating knots with eleven
\cite{Fa}. Lee and Sabloff \cite{LS} extended the computation to fourteen crossings with an
algorithmic search for M\"obius bands in the style of Dunfield and Gong \cite{DG}, together
with the Arf invariant obstruction of Yasuhara \cite{Ya} and a linking form obstruction of
Murakami and Yasuhara \cite{MY}. KnotInfo incorporates these data. In its version
2026.10.5, \ntotal{} knots with eleven to thirteen crossings have an interval $[a,b]$ with
$a<b$ as the entry for $\gamma_4$.

\begin{theorem}\label{thm:gamma}
The value $\gamma_4(K)$ is determined for \nsettled{} of these knots: \nsettledone{} have
$\gamma_4(K)=1$, \nsettledtwo{} have $\gamma_4(K)=2$ and \nsettledthree{} have
$\gamma_4(K)=3$. For \nsettledpinch{} of them a pinch move brings the upper end of the
KnotInfo interval down to its lower end; for \nsettledlower{} the lower bound of
Theorem~\ref{thm:jk} raises the lower end to the upper end; for \nsettledboth{} the two
meet in between. For \nnarrowed{} further knots the interval is narrowed.
\end{theorem}

The knots are listed in Appendix~\ref{app:glist}, and Table~\ref{tab:counts} gives the
counts by crossing number and by KnotInfo interval. Part~I (Sections~\ref{sec:pseudo}
to~\ref{sec:bopen}) proves Theorem~\ref{thm:bridge}, Part~II (Sections~\ref{sec:pinch}
to~\ref{sec:gopen}) proves Theorem~\ref{thm:gamma}. Both parts end with what the methods do
not reach.

\part*{Part I. The bridge index}

\section{A lower bound from pseudo-reflections}\label{sec:pseudo}

A \emph{pseudo-reflection} of a vector space $V$ is an element $g\in\GL(V)$ with
$\rk(g-1)=1$. It fixes the hyperplane $\ker(g-1)$ pointwise.

\begin{lemma}\label{lem:fixed}
Let $V$ be a vector space of dimension $n$ over a field and let $\rho\colon G\to\GL(V)$ be
a homomorphism such that $\rho(m)$ is a pseudo-reflection. Let $V^{G}$ be the subspace of
vectors fixed by $\rho(G)$. Then
\[
\mu(K)\ \ge\ n-\dim V^{G}.
\]
\end{lemma}

\begin{proof}
Let $m_1,\dots,m_k$ be meridians generating $G$, with $m_i=g_im^{\epsilon_i}g_i^{-1}$ and
$\epsilon_i=\pm1$. Since $\rho(m)^{-1}-1=-\rho(m)^{-1}(\rho(m)-1)$, each $\rho(m_i)$
is a pseudo-reflection, and $H_i=\ker(\rho(m_i)-1)$ is a hyperplane. The intersection
$H_1\cap\dots\cap H_k$ has dimension at least $n-k$ and is fixed by every $\rho(m_i)$,
hence by $\rho(G)$. So $n-k\le\dim V^G$.
\end{proof}

\begin{corollary}\label{cor:four}
If $\dim V=4$, $\rho(m)$ is a pseudo-reflection and $V^G=0$, then $b(K)\ge\mu(K)\ge4$.
\end{corollary}

In a diagram with arcs $x_1,\dots,x_r$, the Wirtinger presentation makes such a $\rho$ an
assignment of matrices $X_i=\rho(x_i)$ to the arcs that satisfies
$X_c=X_o^{\pm1}X_aX_o^{\mp1}$ at each crossing, where $x_o$ is the overarc and $x_a,x_c$
are the underarcs. Since the $X_i$ generate $\rho(G)$,
\[
V^G=\bigcap_{i=1}^{r}\ker(X_i-1).
\]
Both conditions are finite computations once the $X_i$ are written down.

\subsection*{Transvections}
Let $\omega$ be a nondegenerate alternating form on $V=\F_q^4$. For $v\ne0$ the
transvection $t_v(x)=x+\omega(x,v)\,v$ lies in $\Sp(V,\omega)$, $t_v-1$ has rank one,
$t_{-v}=t_v$, and $g\,t_v\,g^{-1}=t_{gv}$ for $g\in\Sp(V,\omega)$. A representation with
$X_i=t_{v_i}$ is therefore the same thing as a coloring of the arcs by vectors with
\[
v_c=\pm\,t_{v_o}^{\pm1}(v_a)
\]
at every crossing, and $\ker(X_i-1)=v_i^{\perp}$. So
$V^G=(\operatorname{span}\{v_i\})^{\perp}$, which is $0$ exactly when the $v_i$ span $V$.

\subsection*{Reflections}
Let $B$ be a nondegenerate symmetric bilinear form on $\F_q^4$, $q$ odd. For $B(v,v)\ne0$
the reflection $r_v(x)=x-2\frac{B(x,v)}{B(v,v)}\,v$ satisfies $g\,r_v\,g^{-1}=r_{gv}$ for
$g\in\O(B)$, and again $V^G=(\operatorname{span}\{v_i\})^{\perp}$.

\section{The certificates for the bridge index}\label{sec:bcert}

For each knot of Appendix~\ref{app:bclosed} the file \texttt{data/certificates.json}
contains:
\begin{itemize}
\item the PD code of $K$ from KnotInfo, with edges numbered $1,\dots,2n$ along the
orientation and $[a,b,c,d]$ the edges at a crossing counterclockwise from the incoming
underedge $a$;
\item a prime $q$ and a $4\times4$ matrix over $\F_q$ for every edge;
\item four edges on distinct arcs, the seeds of a Wirtinger coloring.
\end{itemize}
The program \texttt{verify/verify.py} uses only the Python standard library. For every
knot it checks:
\begin{enumerate}
\item the matrices are constant along each arc, where an arc is a maximal run of edges
between two undercrossings;
\item at each crossing $[a,b,c,d]$ the relation $X_c=X_b^{-1}X_aX_b$ holds if the
overstrand runs from $b$ to $d$, and $X_c=X_bX_aX_b^{-1}$ if it runs from $d$ to $b$
(changing this convention at every crossing at once replaces $K$ by its mirror image, which
has the same $b$ and $\mu$);
\item each $X_i-1$ has rank $1$, and the matrices $X_i-1$ stacked together have rank $4$,
so that $V^G=0$;
\item starting from the four seed arcs, the coloring moves of \cite{BKVV} (if the overarc
and one underarc at a crossing are colored, color the other underarc) color every arc, so
$\omega(D)\le4$.
\end{enumerate}
By Corollary~\ref{cor:four} and \eqref{eq:upper}, items (1)--(4) prove $b(K)=\mu(K)=4$.
With the option \texttt{--knotinfo} the program also checks that each PD code is the one
in KnotInfo and that KnotInfo lists the bridge index as $[3,4]$. All \bclosed{}
certificates are checked in a few seconds. The program rejects altered certificates: a
changed matrix entry, a seed set on three arcs, one crossing with its overstrand reversed,
or identity matrices.

\begin{remark}
The certificates were found by a search, which plays no part in the proof. If the four
seed vectors do not span $V$, every arc vector lies in their span, so a certificate must
have seeds forming a basis. A coloring is then determined by the Gram matrix of the seeds:
six entries of an alternating matrix, or six off-diagonal entries of a symmetric matrix
with a fixed diagonal. The search runs over all such matrices and propagates the colors
from the seeds through the crossings. The Wirtinger relations are then checked at the
remaining crossings.
\end{remark}

\section{The knots whose bridge index remains open}\label{sec:bopen}

For the other \bopen{} knots with entry $[3,4]$ (Appendix~\ref{app:bopen}) we found
neither bound. The exhaustive searches described above found no representation satisfying
Corollary~\ref{cor:four} with meridians sent to transvections of $\Sp_4(\F_q)$ for
$q\in\{3,4,5,7,8,9,11\}$, to reflections of $\O_4(\F_q)$ for $q\in\{3,5,7\}$ (both
isometry types), or to pseudo-reflections of $\GL_4(\F_2)$ and $\GL_4(\F_3)$ with
eigenvalue $1$ or $-1$. Random Reidemeister moves (SnapPy \cite{SnapPy}) produced no
diagram of Wirtinger number $3$; 300 diagrams were tried for each knot. Abelian colorings
give nothing either: for every prime $p<110$ and every irreducible factor $g$ of the
Alexander polynomial modulo $p$, the colorings of the diagram by the Alexander quandle
$\F_p[t]/(g)$, $a\triangleright b=ta+(1-t)b$, form a space of dimension at most $3$, while
dimension $4$ would give $\mu(K)\ge4$. Nor is any of these knots the closure of a
$3$-braid $\Delta^{2d}w$ in Murasugi's normal form with $|d|\le4$ and $w$ either
$\sigma_1^{-p_1}\sigma_2^{q_1}\cdots\sigma_1^{-p_s}\sigma_2^{q_s}$ of length at most
$17-4\max(|d|-1,0)$ or $\sigma_2^m$ with $|m|\le20$ ($53{,}309$ braids, identified by
SnapPy), which would give $b(K)\le3$. A diagram of Wirtinger number $3$ would give
$b(K)=\mu(K)=3$, and a representation as in Corollary~\ref{cor:four} over another field,
or a Coxeter quotient of rank $4$, would give $b(K)=\mu(K)=4$.

\part*{Part II. The nonorientable four-genus}

\section{Pinch moves}\label{sec:pinch}

\begin{definition}
Let $D$ be a diagram of a knot $K$ and $c$ a crossing of $D$. Of the two smoothings of $c$,
one respects the orientation of $K$ and gives a two-component link; the other gives a knot
$K'$. We say $K'$ is obtained from $K$ by a \emph{pinch move} at $c$.
\end{definition}

\begin{lemma}\label{lem:pinch}
If $K'$ is obtained from $K$ by a pinch move, then $\gamma_4(K)\le 1$ when $K'$ is slice,
and $\gamma_4(K)\le\gamma_4(K')+1$ in general.
\end{lemma}

\begin{proof}
The pinch move is a band move, so there is a cobordism $\Sigma\subset S^3\times[0,1]$ from
$K$ to $K'$ formed by $K\times[0,1]$ and one band. Since $K'$ is connected, the band does
not respect the orientation of $K$, and $\Sigma$ is a M\"obius band with an open disk
removed. If $K'$ bounds a disk $\Delta$ in $B^4$, then $\Sigma\cup\Delta$ is a M\"obius band
bounded by $K$, so $\gamma_4(K)\le1$. If $K'$ bounds a surface $F$ with
$b_1(F)=\gamma_4(K')$, then $\Sigma\cup F$ is a nonorientable surface bounded by $K$ with
$b_1=b_1(F)+1$. This is the bound used by Jabuka and Kelly \cite{JK}.
\end{proof}

A knot that is not slice has $\gamma_4\ge1$. For the knots whose KnotInfo interval starts
at $2$, the lower bound is the obstruction to bounding a M\"obius band recorded by KnotInfo
from Lee and Sabloff \cite{LS}: either Yasuhara's condition
$4\operatorname{Arf}(K)-\sigma(K)\in\{0,\pm2\}\pmod 8$ fails \cite{Ya}, or the linking form
obstruction of Murakami and Yasuhara \cite{MY} applies. The further lower bounds are those
of Section~\ref{sec:lower}. Theorem~\ref{thm:gamma} therefore rests on these obstructions,
on the slice status and four-genus values recorded in KnotInfo for the targets, on the
certificates of Section~\ref{sec:gcert} and on the lattice computations of
Section~\ref{sec:lower}.

\section{The pinch certificates and their verification}\label{sec:gcert}

A certificate for a knot $K$ consists of a planar diagram code $D$, a crossing $c$ of $D$, a
choice of smoothing, and a target: either the unknot or a knot $T$ of KnotInfo. It asserts
that
\begin{enumerate}
\item $D$ is a diagram of $K$;
\item the chosen smoothing of $c$ gives a one-component diagram $D'$;
\item $D'$ represents $T$ or its mirror image;
\item the bound $\beta=1$ if $T$ is slice or the unknot, and $\beta=\gamma_4(T)+1$ otherwise,
is smaller than the upper end of the KnotInfo interval of $K$.
\end{enumerate}
Since $\gamma_4$ and sliceness are invariant under mirror image, (1) to (4) and
Lemma~\ref{lem:pinch} give $\gamma_4(K)\le\beta$.

The certificates were found with SnapPy \cite{SnapPy}: for every open knot we applied
every pinch move to the KnotInfo diagram and to up to twenty-five diagrams obtained from it
by random Reidemeister moves, and identified the resulting knots by hyperbolic structure. A
candidate enters the file only if the verifier below accepts it; when it did not, the search
went on with other crossings and other diagrams of the same knot. A second pass used the
knots settled in the first pass as targets; the number of certificates with such a target
is \nchained{}, and the verifier checks them after the certificates of their targets. A last
pass applied every pinch move to twenty diagrams of each knot left open.

The verifier \texttt{crosscap4d/verify/verify.py} does not use SnapPy. It reads the KnotInfo
database and uses Regina \cite{Regina} for the topology. For (1) and (3) it proves that two
diagrams represent the same knot up to mirror image in one of two ways.
\begin{itemize}
\item \emph{Isotopy.} Both diagrams are simplified by Reidemeister moves to the same number
of crossings, and the closures of the two minimal diagrams under Reidemeister moves that add
at most one crossing share a diagram. Diagrams are compared by Regina's canonical
signatures, taken up to reflection, rotation and reversal.
\item \emph{Complements.} The triangulated complements of the two diagrams have
combinatorially isomorphic triangulations, possibly after Pachner moves that add at most one
tetrahedron, or one of the triangulations obtained from forty simplifications of each
complement under fixed seeds of Regina's random engine, together with the triangulations one
Pachner move away from the smallest of them, occurs for both diagrams. Each of these
triangulations is a triangulation of the complement it came from, so a common one shows that
the complements are homeomorphic, and by the theorem of Gordon and Luecke \cite{GL} the
knots agree up to mirror image.
\end{itemize}
Both tests are exact: a positive answer is a proof. Regina's simplification is randomised,
so the verifier makes up to four attempts, each from a fixed seed, and gives each
certificate two minutes. For the unknot the verifier uses Regina's recognition of the solid
torus. Item (2) is checked by following the strands of $D'$, and (4) by reading KnotInfo.

The verifier accepts all \ncerts{} certificates. It rejects certificates corrupted in four
ways, on a sample of twenty: the other smoothing of the crossing, the next crossing, a wrong
target, and a bound lowered by one (Table~\ref{tab:selftest}). The seven modified
certificates with the next crossing that it accepts are valid certificates: that crossing
gives the same target.

\begin{table}[htbp]
\centering
\selftesttable
\caption{Corrupted pinch certificates rejected by the verifier.}
\label{tab:selftest}
\end{table}

\section{Lower bounds from definite lattices}\label{sec:lower}

Let $D$ be a diagram of a knot $K$ with a checkerboard colouring. Following Gordon and
Litherland \cite{GoL}, give each crossing $c$ the weight $\eta(c)=\pm1$ according to the
position of the shaded regions at $c$, let $X_0,\dots,X_n$ be the unshaded regions, and let
$G$ be the $n\times n$ matrix with $g_{ij}=-\sum\eta(c)$ over the crossings meeting $X_i$
and $X_j$ for $i\ne j$, and $g_{ii}=-\sum_{k\ne i}g_{ik}$ with $k$ running from $0$ to $n$.
This is the Goeritz matrix; $\det G=\pm\det K$, and with $\mu=\sum\eta(c)$ over the crossings
of type~II for the shaded surface,
\begin{equation}\label{eq:gl}
\sigma(K)=\operatorname{sign}(G)-\mu .
\end{equation}
The form $G$ is the intersection form of the double cover of $B^4$ branched along the shaded
surface pushed into $B^4$ \cite{GoL}. If $D$ is reduced and alternating, the forms $G_+$ and
$G_-$ of the two colourings are positive and negative definite.

\begin{theorem}[Jabuka and Kelly {\cite[Theorems 2.10--2.12]{JK}}]\label{thm:jk}
Let $K$ be a knot with a reduced alternating diagram, let $P_+=G_+$ and $P_-=-G_-$, of ranks
$r_\pm$, and let $s\equiv\sigma(K)+4\operatorname{Arf}(K)\pmod 8$. Call $\ell\in\mathbb N$
admissible if $\ell\mid\det K$ and $\det K/\ell$ is a square.
\begin{enumerate}
\item If $s=2$ and $P_-\oplus[\ell]$ embeds in $\mathbb Z^{r_-+1}$ for no admissible
$\ell$, or $s=6$ and $P_+\oplus[\ell]$ embeds in $\mathbb Z^{r_++1}$ for no admissible
$\ell$, then $\gamma_4(K)\ge2$.
\item If $s=0$ and neither $P_+\oplus[\ell]$ embeds in $\mathbb Z^{r_++1}$ nor
$P_-\oplus[\ell]$ in $\mathbb Z^{r_-+1}$, for any admissible $\ell$, then $\gamma_4(K)\ge2$.
\item If $s=4$ and neither $P_+$ embeds in $\mathbb Z^{r_++2}$ nor $P_-$ in
$\mathbb Z^{r_-+2}$, then $\gamma_4(K)\ge3$.
\end{enumerate}
Here $\mathbb Z^N$ carries the standard positive definite form.
\end{theorem}

We recall the argument. If $K$ bounds a M\"obius band $\Sigma\subset B^4$, the double
branched cover $W(\Sigma)$ has $b_1=0$ \cite{JK} and $b_2=1$, its form is $[\pm\ell]$ with
$\ell$ admissible by a theorem of Gilmer and of Owens and Strle \cite[Corollary 2.8]{JK},
and the congruence of Gilmer and Livingston \cite{GiL} fixes the sign of $\ell$ by $s$.
Gluing $W(\Sigma)$ with reversed orientation to the double cover bounded by the shaded
surface whose Goeritz form has the opposite sign gives a closed definite smooth
four-manifold whose form contains $G\oplus[\mp\ell]$ in full rank. By Donaldson's theorem
\cite{Do} that form is diagonal; changing the sign of the whole form gives (1) and (2). In
(3) a surface with $b_1=2$ would give a definite $W(\Sigma)$ of rank two, of either sign,
and the same gluing embeds $P_+$ or $P_-$ in a diagonal lattice of rank two more. In (1) the
statement of \cite{JK} names the target lattice as $\varepsilon$-definite; the argument
above, which is theirs, produces the lattice of the sign of $G$, and it is in this form that
we apply it.

Since $\gamma_4$ is invariant under mirror image, and mirroring changes the sign of every
weight $\eta(c)$, of $G$, of $\mu$ and of $\sigma$, it does not matter which of the two
weight conventions is called positive, as long as $\sigma$ is computed from \eqref{eq:gl}
with the same weights.

\begin{lemma}\label{lem:complement}
Let $P$ be a positive definite lattice of rank $r$ and $\ell\in\mathbb N$. Then
$P\oplus[\ell]$ embeds in $\mathbb Z^{r+1}$ if and only if there are an embedding
$\phi\colon P\to\mathbb Z^{r+1}$ and an integer $k\ge1$ with $\ell=k^2|w|^2$, where $w$ is a
primitive vector spanning the orthogonal complement of $\phi(P)$.
\end{lemma}

\begin{proof}
The orthogonal complement of a sublattice of rank $r$ in $\mathbb Z^{r+1}$ is $\mathbb Z w$
for a primitive $w$, and a vector of norm $\ell$ orthogonal to $\phi(P)$ is $kw$ with
$k^2|w|^2=\ell$.
\end{proof}

\begin{lemma}\label{lem:search}
Let $P$ be a positive definite Gram matrix of rank $r$ and $N\ge r$. Every embedding of $P$
in $\mathbb Z^N$ is carried by a signed permutation of the coordinates to one found by the
following search. The basis vectors $e_1,\dots,e_r$ are placed in turn; if the vectors
already placed are supported on the first $t$ coordinates, the next vector is
$(x,y,0)$ with $x\in\mathbb Z^t$ arbitrary, subject to the norm and the inner products with
the vectors already placed, and $y$ a nonincreasing sequence of positive integers.
\end{lemma}

\begin{proof}
By induction on the number of vectors placed. Given an embedding $v_1,\dots,v_r$ whose first
$i-1$ vectors are of the stated form, supported on the first $t$ coordinates, the signed
permutations of the coordinates beyond $t$ fix $v_1,\dots,v_{i-1}$; one of them makes the
entries of $v_i$ beyond $t$ nonnegative and nonincreasing, with its nonzero entries first,
and it does not change the shape of the vectors already placed.
\end{proof}

We take the KnotInfo diagram of each open alternating knot, compute both Goeritz forms and
$\sigma$ from \eqref{eq:gl}, check that the two colourings give the same $\sigma$, that
$|\sigma|$ agrees with KnotInfo and that the forms are definite, and run the search of
Lemma~\ref{lem:search} to the end: for (1) and (2) it lists the norms $|w|^2$ of
Lemma~\ref{lem:complement} over all embeddings, for (3) it decides whether an embedding
exists. Of the \ntotal{} knots, the hypotheses of Theorem~\ref{thm:jk} hold and raise the
KnotInfo lower end for \nlower{}: \nlowertwo{} reach $\gamma_4\ge2$ (\nlowerten{} by (1)
and \nlowertwelve{} by (2)) and \nlowerthree{} reach $\gamma_4\ge3$ by (3). The data are in
\texttt{crosscap4d/data/lower\_bounds.json}.

The program \texttt{crosscap4d/lower/verify\_lower.py} recomputes every entry with the basis
order of the search shuffled under a fixed seed, and runs three tests of the conventions. On
the alternating knots with at most eleven crossings whose $\gamma_4$ KnotInfo records, the
bound never exceeds the recorded value, while the same program with the sign of $\sigma$
reversed contradicts the recorded value for more than two hundred of them. On the
alternating knots that carry a pinch certificate, the bound never exceeds the upper bound of
that certificate.

\begin{table}[htbp]
\centering
\countstable
\caption{Open entries of $\gamma_4$ in KnotInfo 2026.10.5 and what the certificates give.
``Settled'': the value is now known; ``lowered'': the upper bound is lower.}
\label{tab:counts}
\end{table}

\section{The knots whose nonorientable four-genus remains open}\label{sec:gopen}

The remaining \nopen{} knots, with their narrowed intervals, are listed in the
file \texttt{still\_open.json} of the folder \texttt{crosscap4d/data}. For a knot with interval $[1,2]$ that is not
settled, no pinch move we tried gives a slice knot, and Theorem~\ref{thm:jk} does not apply
or does not obstruct. Theorem~\ref{thm:jk} needs a definite diagram, so it says nothing
about the nonalternating knots.

The bound $|\upsilon(K)-\sigma(K)/2|\le\gamma_4(K)$ of Ozsv\'ath, Stipsicz and Szab\'o
\cite{OSS}, with $\upsilon=\Upsilon_K(1)$, does not raise any of these lower ends where we
could evaluate it. For the open knots with Hom's invariant $\varepsilon(K)=0$ the complex
$CFK^\infty(K)$ is $\varepsilon$-equivalent to that of the unknot \cite{Hom}, so
$\upsilon(K)=0$ and the bound is $|\sigma(K)|/2$, which never exceeds the lower end of the
KnotInfo interval. For the others, the knot Floer complexes we could compute are complexes over
$\F[U,V]/(UV)$, which do not determine $\Upsilon_K(1)$.

\begin{remark}
The \nclosedone{} M\"obius bands found here are not in the data of Lee and Sabloff
\cite{LS}, whose search follows Dunfield and Gong \cite{DG} and runs over a finite family of
bands chosen in advance. The knots settled at $\gamma_4=2$ were outside their aim, which was
M\"obius bands; for them the target of the pinch move is a knot known to bound a M\"obius
band but not a disk.
\end{remark}

\section*{Data availability}
The certificates, the verifiers and the search programs for both parts are in the
repository \url{https://github.com/creelie/conjecture-to-choose}, release 1.3.0,
\url{https://github.com/creelie/conjecture-to-choose/releases/tag/v1.3.0}: Part~I in the
folders \texttt{data}, \texttt{verify} and \texttt{search}, Part~II in the folder
\texttt{crosscap4d}.

\section*{Use of AI tools}
The searches, the verification programs and drafts of the text were produced with Claude
Code (Anthropic) under the direction of the author, who is responsible for the
content.

\section*{Funding and competing interests}
The author received no funding for this work and has no competing interests.

\begin{thebibliography}{99}
\bibitem{BKM} R.~Blair, A.~Kjuchukova and N.~Morrison, \emph{Coxeter quotients of knot
groups through 16 crossings}, Exp. Math. 34 (2025), no.~1, 27--37; arXiv:2208.09032.
\bibitem{BKP} R.~Blair, A.~Kjuchukova and E.~Pfaff, \emph{The plain sphere number of a
link}, arXiv:2504.10517.
\bibitem{BKVV} R.~Blair, A.~Kjuchukova, R.~Velazquez and P.~Villanueva, \emph{Wirtinger
systems of generators of knot groups}, Comm. Anal. Geom. 28 (2020), no.~2, 243--262;
arXiv:1705.03108.
\bibitem{Regina} B.~A.~Burton, R.~Budney, W.~Pettersson et al., \emph{Regina: software for
low-dimensional topology}, regina-normal.github.io, version 7.4.
\bibitem{SnapPy} M.~Culler, N.~M.~Dunfield, M.~Goerner and J.~R.~Weeks, \emph{SnapPy, a
computer program for studying the geometry and topology of 3-manifolds}, version 3.3.2,
\url{https://snappy.computop.org}.
\bibitem{Do} S.~K.~Donaldson, \emph{The orientation of Yang--Mills moduli spaces and
4-manifold topology}, J. Differential Geom. 26 (1987), 397--428.
\bibitem{DG} N.~M.~Dunfield and S.~Gong, \emph{Ribbon concordances and slice obstructions:
experiments and examples}, arXiv:2512.21825.
\bibitem{Fa} M.~Fairchild, \emph{The non-orientable 4-genus of 11 crossing non-alternating
knots}, J. Knot Theory Ramifications 34 (2025), no.~1.
\bibitem{Gh} N.~Ghanbarian, \emph{The non-orientable 4-genus for knots with 10 crossings},
J. Knot Theory Ramifications 31 (2022), no.~5, 2250034.
\bibitem{GiL} P.~M.~Gilmer and C.~Livingston, \emph{The nonorientable 4-genus of knots},
J. Lond. Math. Soc. (2) 84 (2011), 559--577.
\bibitem{GoL} C.~McA.~Gordon and R.~A.~Litherland, \emph{On the signature of a link},
Invent. Math. 47 (1978), 53--69.
\bibitem{GL} C.~McA.~Gordon and J.~Luecke, \emph{Knots are determined by their complements},
J. Amer. Math. Soc. 2 (1989), 371--415.
\bibitem{Hom} J.~Hom, \emph{A survey on Heegaard Floer homology and concordance}, J. Knot
Theory Ramifications 26 (2017), no.~2, 1740015.
\bibitem{JK} S.~Jabuka and T.~Kelly, \emph{The nonorientable 4-genus for knots with 8 or 9
crossings}, Algebr. Geom. Topol. 18 (2018), 1823--1856.
\bibitem{Kirby} R.~Kirby (ed.), \emph{Problems in low-dimensional topology}, in Geometric
Topology (Athens, GA, 1993), AMS/IP Stud. Adv. Math. 2.2, Amer. Math. Soc., 1997, 35--473.
\bibitem{LS} D.~Lee and J.~M.~Sabloff, \emph{An algorithmic search for knots bounding
M\"obius bands in $B^4$}, arXiv:2607.23582.
\bibitem{KnotInfo} C.~Livingston and A.~H.~Moore, \emph{KnotInfo: Table of Knot Invariants},
\url{https://knotinfo.org}, version 2026.10.5 (the \texttt{database\_knotinfo} package).
\bibitem{MY} H.~Murakami and A.~Yasuhara, \emph{Four-genus and four-dimensional clasp number
of a knot}, Proc. Amer. Math. Soc. 128 (2000), 3693--3699.
\bibitem{OSS} P.~S.~Ozsv\'ath, A.~I.~Stipsicz and Z.~Szab\'o, \emph{Unoriented knot Floer
homology and the unoriented four-ball genus}, Int. Math. Res. Not. IMRN 2017, no.~17,
5137--5181.
\bibitem{Ya} A.~Yasuhara, \emph{Connecting lemmas and representing homology classes of simply
connected 4-manifolds}, Tokyo J. Math. 19 (1996), 245--261.
\end{thebibliography}

\appendix
\section{The knots of Theorem~\ref{thm:bridge}}\label{app:bclosed}
{\small\raggedright
\bclosedlist.\par}

\section{The knots with bridge index entry \texorpdfstring{$[3,4]$}{[3,4]} that remain open}\label{app:bopen}
{\small\raggedright
\bopenlist.\par}

\section{The knots of Theorem~\ref{thm:gamma}}\label{app:glist}

{\footnotesize
\noindent\textbf{$\gamma_4=1$ (\nsettledone{} knots).}
\begin{multicols}{6}\raggedright
\listone
\end{multicols}

\noindent\textbf{$\gamma_4=2$ (\nsettledtwo{} knots).}
\begin{multicols}{6}\raggedright
\listtwo
\end{multicols}

\noindent\textbf{$\gamma_4=3$ (\nsettledthree{} knots).}
\begin{multicols}{6}\raggedright
\listthree
\end{multicols}
}

\end{document}
